\chapter{软件架构设计与模块划分}

\paragraph*{学习目标}

通过本章学习，学生应能够：
\begin{enumerate}
\item
  会把“高性能”“高可靠”一类形容词需求改写为含刺激、环境、响应与度量的质量属性场景，并给出验证手段
\item
  运用UML建模语言描述软件架构，能够绘制用例图、类图、活动图、状态图和时序图
\item
  理解并掌握智慧水利平台模块划分原则和详细设计方法，能够进行合理的系统分解和接口设计
\item
  能够形成系统架构图、模块职责、关键接口和架构决策等设计文档
\item
  了解软件架构评估方法和质量保证技术，能够对设计方案进行评价和优化
\end{enumerate}

\paragraph*{引言}

软件架构设计是智慧水利平台开发的核心环节，承接需求分析的成果，支撑后续的编码实现。在复杂的智慧水利系统中，合理的模块划分和详细设计影响系统的功能实现、性能、可维护性、可扩展性和可靠性。

智慧水利平台涉及数据采集、传输、存储、处理、分析、展示等多个技术领域，并需要集成水文监测、工程管理、应急调度、决策支持等多种业务功能。这种技术与业务的复杂性要求我们必须采用科学的设计方法，通过合理的模块化分解，实现复杂系统的有序构建。

\section{软件架构概述}

\subsection{软件架构的定义与重要性}

软件架构是一个系统的“结构或结构集合，包括软件构件、构件间的关系以及构件与环境的关系”。它定义了系统的整体框架，决定了系统如何实现功能、处理数据以及与外部交互。架构设计是软件工程的核心活动之一，对系统的质量、可维护性和扩展性有重要影响\cite{ISO42010,clements2010,richards2020,shaw1996}。

\paragraph{架构对系统质量的影响}

软件架构是系统的基石，影响系统的功能实现，也影响性能、可靠性、安全性、可维护性和可移植性这几类非功能性需求能否达成。质量模型可用于组织和评价这些质量属性\cite{ISO25010,GB2006}。

\paragraph{性能}智慧水利平台面临实时数据处理的严峻挑战。大量监测设备（渗压计、GNSS 监测站、雨量站等）持续产生数据，平台需要按表\ref{tab:performance-baseline}分别控制查询、告警生成和推送送达的 P95。是否需要分布式部署、负载均衡和异步处理，取决于测点规模、并发量和可用性目标：案例水库这类单工程平台，一个组织良好的模块化单体加消息队列即可达标；只有当汛期峰值或多工程接入超出单机容量时，才值得为分布式付出运维复杂度的代价——这正是后文架构权衡要展开的判断。例如，通过分离表示层和业务逻辑层，监测数据采集、处理和展示可以独立扩展；利用消息队列（如 Kafka）实现异步处理，使数据流通过中间件而非直接调用，避免任何单一节点成为性能瓶颈。在分布式架构下，多个服务器节点可以并行处理汛期高峰期的数据量激增。

\paragraph{可靠性}水利工程的安全性与生命财产密切相关，监测系统对可靠性的要求显著高于一般信息系统；工程上不存在“任何条件下都不中断”的系统，可行的目标是明确可用性指标（如全年可用率、允许的单次中断时长），并为超出设计范围的故障准备降级和恢复路径。合理的架构通过冗余设计、故障转移和自动恢复机制来降低单点故障的风险。典型的做法包括：（1）采用主从架构或分布式集群，监测数据同时写入多个副本，当某个节点故障时其他节点可立即接管；（2）部署健康检查和自动重启机制，让异常进程的恢复时间从人工介入的小时级缩短到自动化的分钟级，恢复目标应写入可用性指标而不是凭直觉承诺；（3）实现完整的错误处理和事务管理，保证数据一致性，确保监测值、告警记录和处置指令都准确可追溯；（4）建立冗余的网络链路和通信协议（如 MQTT 重连和消息持久化），应对网络波动。这些架构层的容错设计直接关系到防洪应急、工程安全等关键业务的连续性。

表\ref{tab:performance-baseline}给出本章案例水库用于架构评审的假定性能基线；全书工程参数以第8章8.1节参数表为唯一来源，真实项目再按并发量、网络条件和业务时限确认。
\begin{table}[htbp]
\centering
\caption{案例水库性能基线}\label{tab:performance-baseline}
\begin{tabular}{p{4.2cm}p{3.2cm}p{6.0cm}}
\toprule
操作 & 指标 & 度量边界 \\
\midrule
交互查询 & P95 $\leq 1$ s & 从 API 收到查询到返回页面所需数据 \\
告警生成 & P95 $\leq 5$ s & 数据入站到产生告警事件 \\
告警推送送达 & P95 $\leq 10$ s & 事件生成到值班终端收到通知 \\
\bottomrule
\end{tabular}
\end{table}

\begin{table}[htbp]
\centering
\caption{架构质量属性与水利场景对应关系}\label{tab:quality-attributes}
\begin{tabular}{p{1.8cm}p{5.0cm}p{5.5cm}}
\toprule
质量属性 & 架构决策与实现 & 水利场景应用 \\
\midrule
性能 & 异步处理、缓存、批流协同、水平扩展 & 汛期高峰数据接入与告警查询保持稳定响应 \\
可靠性 & 冗余部署、故障转移、消息持久化、恢复演练 & 设备或节点故障时持续采集，告警与处置记录不丢失 \\
安全性 & 加密通信、基于角色的访问控制、审计日志 & 监测数据防篡改，关键操作授权并留痕 \\
可维护性 & 监测接入、处理与预警模块分离，接口规范化 & 更新设备驱动不影响预警模块，故障定位边界清晰 \\
扩展性 & 分层架构、插件式适配、松耦合服务 & 接入新设备和分析模型，扩展到河道、灌区等场景 \\
\bottomrule
\end{tabular}
\end{table}

表\ref{tab:quality-attributes}列出的五类属性是本案例的工程化关注项。ISO/IEC 25010 的八项顶层质量特性为功能适合性、性能效率、兼容性、易用性（usability）、可靠性、安全性、可维护性和可移植性；其中表中的性能、可靠性、安全性、可维护性分别对应性能效率、可靠性、安全性和可维护性，扩展性是面向演进的工程指标，可由可维护性、兼容性和可移植性共同支撑。架构决策必须权衡它们的优先级：对于关系生命安全的防洪系统，可靠性和安全性优先；对于汛期数据洪峰，性能与扩展性更受关注；可维护性则保障系统长期演进。决策时还应记录需求、备选方案、成本、风险和选择理由，使质量目标能够被测试和复核。

质量属性只有写成\textbf{场景}才能被设计和检验。一条质量属性场景包含六要素：刺激源、刺激、制品、环境、响应、响应度量。“系统要高可靠”不可设计；写成场景就变成：“汛期峰值期间（环境），采集网关每秒推送500条观测（刺激源与刺激）到接入服务（制品），系统应在不丢数据的前提下完成质量检查并入库（响应），端到端时延 P95 不超过5秒、丢失率为零（度量）。”案例水库的核心场景至少应包括：性能场景（对应 REQ-MON-03，度量取表\ref{tab:performance-baseline}）、可靠性场景（单节点故障后曲线展示降级为缓存快照，恢复时间不超过5分钟）、安全场景（未授权用户请求修改阈值，系统拒绝并产生审计事件，对应 REQ-MON-04）、可维护性场景（新增一类传感器只改动协议适配层，两人日内完成）。本章后续的架构决策、案例评审与ATAM演练，都以这组场景为公共输入——场景写不出来，说明需求还停留在形容词阶段，应退回第2章的方法补课。

表\ref{tab:ch03-quality-scenarios}把这四条核心场景按六要素展开，右列同时给出验证手段——场景与验证手段必须成对出现，这是它区别于愿望清单的地方。

\begin{table}[htbp]
\centering
\caption{案例水库的核心质量属性场景}
\label{tab:ch03-quality-scenarios}
\begin{tabular}{p{1.6cm}p{4.6cm}p{3.6cm}p{3.2cm}}
\toprule
属性 & 刺激与环境 & 响应与度量 & 验证手段 \\
\midrule
性能 & 汛期峰值，网关每秒推送500条观测到接入服务 & 完成质检并入库，端到端P95不超过5秒，零丢失 & 压测脚本回放峰值样本 \\
可靠性 & 常规运行中单个服务节点故障 & 曲线降级为带时间标注的缓存快照，5分钟内恢复 & 停机演练与恢复计时 \\
安全 & 未授权用户请求修改预警阈值 & 拒绝操作并产生含主体与时间的审计事件 & 越权用例自动化测试 \\
可维护性 & 新增一类光纤渗压传感器 & 只改动协议适配层，两人日内接入完成 & 适配层改动的代码评审 \\
\bottomrule
\end{tabular}
\end{table}


\subsection{架构描述方法}

软件架构的描述方法是将抽象的系统结构转化为可视化、可理解的模型，以指导开发与协作\cite{clements2010}。根据ISO/IEC/IEEE 42010标准，架构描述应当包含多个视图，每个视图从特定的利益相关者角度展现系统架构\cite{ISO42010}。UML（统一建模语言）提供了多种视图，帮助工程师从不同角度描述架构；而架构决策的核心要素则确保设计满足系统的核心需求。

\subsubsection{UML视图}
统一建模语言（Unified Modeling Language，UML）是一种用于可视化、详述、构造和文档化软件系统以及其他非软件系统的建模语言。UML通过多种视图全面描述系统的架构，每个视图都从不同的角度和关注点展现系统的特定方面，从而为开发者、分析师、测试人员和用户等不同角色提供了理解系统的统一框架。以下是主要视图及其作用：

\paragraph{用例模型}

用例图（Use Case Diagram）是UML中用来描述系统功能需求的重要视图。它说明系统与外部参与者之间的交互方式，明确系统的边界以及用户与系统之间的行为场景。用例图通过用例、参与者和它们之间的关系，帮助开发团队理解系统的功能要求和使用场景。

\paragraph{用例}用例代表系统所提供的一个功能或服务，它是系统行为的抽象描述。每个用例都应当满足某种业务需求或用户需求。在用例图中，用例通常用椭圆形表示，标注出其名称。用例能够清晰展示系统的业务功能或处理过程，比如在智慧水利平台中，常见的用例有“查看监测大屏”、“查询测站数据”、“确认告警事件”等。

\paragraph{参与者}参与者是指与系统交互的外部实体，可以是人、外部系统或其他软件组件。参与者通过与用例的交互来触发系统的行为。在用例图中，参与者用带有小人图标的角色表示，如“值班员”“专业分析员”“审批人”“运维员”或外部监测设备平台。

\paragraph{关系}

关联关系（Association）：表示参与者和用例之间的交互，通常表现为一条简单的实线。

包含关系（Include）：表示某个用例包含了其他用例的功能，即一个用例的一部分。它表示模块化的功能结构。

扩展关系（Extend）：表示用例在特定条件下可能会扩展其功能，增强或变更基本用例的行为。

泛化关系（Generalization）：表示参与者或用例之间的继承关系，即某个用例或参与者是另一个用例或参与者的特化版本。

例如，在水利工程安全监测平台中，值班员、专业分析员、审批人和运维员会通过监测大屏、Web 页面或移动终端与系统交互，完成监测查询、预警确认、研判审批、设备维护和报表分析等操作。用例图可以描述这些交互场景，帮助明确平台需要提供的核心业务能力以及它们之间的关系。

图\ref{fig:uml-usecase}给出一个最小用例图。系统边界内放置业务能力，边界外放置参与者；连线表示参与关系，而不是调用顺序。

\begin{figure}[H]
\centering
\begin{tikzpicture}[>=Stealth, every node/.style={font=\small}]
  \draw[thick, draw=Primary] (-1.8,-2.0) rectangle (5.2,2.0);
  \node[font=\small\bfseries, text=PrimaryDark] at (1.7,1.62) {水利工程安全监测平台};
  \node[draw, circle, minimum size=5mm] (operator) at (-3.4,0.8) {};
  \draw (operator.south) -- ++(0,-0.6) coordinate (obody);
  \draw ($(operator.south)+(0,-0.2)$) -- ++(-0.35,-0.25);
  \draw ($(operator.south)+(0,-0.2)$) -- ++(0.35,-0.25);
  \draw (obody) -- ++(-0.3,-0.45); \draw (obody) -- ++(0.3,-0.45);
  \node[below=1.15cm of operator] {值班员};
  \node[draw=Primary, ellipse, fill=PrimaryLight, minimum width=2.4cm] (view) at (0,0.9) {查看监测态势};
  \node[draw=Primary, ellipse, fill=PrimaryLight, minimum width=2.4cm] (ack) at (0,-0.8) {确认告警事件};
  \node[draw=Primary, ellipse, fill=PrimaryLight, minimum width=2.4cm] (rule) at (3.0,0.2) {维护预警规则};
  \node[draw=Primary, ellipse, fill=PrimaryLight, minimum width=2.4cm] (log) at (2.0,-0.8) {记录处置日志};
  \node[draw, circle, minimum size=5mm] (managerHead) at (6.4,0.8) {};
  \draw (managerHead.south) -- ++(0,-0.6) coordinate (managerBody);
  \draw ($(managerHead.south)+(0,-0.2)$) -- ++(-0.35,-0.25);
  \draw ($(managerHead.south)+(0,-0.2)$) -- ++(0.35,-0.25);
  \draw (managerBody) -- ++(-0.3,-0.45); \draw (managerBody) -- ++(0.3,-0.45);
  \node[below=1.15cm of managerHead] {专业分析员};
  \draw (operator.east) -- (view.west); \draw (operator.east) -- (ack.west);
  \draw (managerHead.west) -- (rule.east);
  \draw[dashed,->] (ack) -- node[above, font=\scriptsize, fill=white, inner sep=1pt] {«include»} (log);
  \node[draw=Primary, ellipse, fill=white, minimum width=2.4cm, font=\scriptsize] (hist) at (3.2,1.35) {查看历史同期对比};
  \draw[dashed,->] (hist) -- node[right, font=\scriptsize, fill=white, inner sep=1pt] {«extend»} (view);
\end{tikzpicture}
\caption{水利工程安全监测平台用例图示例}
\label{fig:uml-usecase}
\end{figure}

\paragraph{设计模型}

设计模型是对系统内部结构的详细描述，说明系统各模块如何实现业务需求。以下是几种常见的设计模型视图：

\paragraph{（1）类图（Class Diagram）}类图是UML中最常用的设计图之一，用于展示系统中的类及其属性、操作（方法）以及类与类之间的关系。类图为面向对象设计提供了基础，帮助开发者理解系统的数据结构、类之间的相互作用和依赖关系。通过类图，可以看到类之间的继承、关联、依赖等关系。

\paragraph{类}类是对象的蓝图，包含属性（变量）和方法（操作）。类图中的类通常表示为矩形框，框内标注类的名称、属性和方法。

\paragraph{关系}

泛化（继承）关系（Generalization）：表示一个类继承了另一个类的属性和方法，使用实线加空心三角箭头，三角形指向父类；UML 规范术语为泛化。

关联关系（Association）：表示类之间的联系，通常用一条实线表示。

依赖关系（Dependency）：表示一个类依赖于另一个类来完成某个操作，通常用带箭头的虚线表示。

聚合与组合关系（Aggregation and Composition）：表示类与类之间的整体与部分关系，聚合使用整体端的空心菱形，组合使用整体端的实心菱形；菱形位置比“强弱”文字更能明确整体端归属。

实现关系表示类对接口的实现，使用虚线加空心三角；依赖关系使用虚线加开放箭头；关联关系使用实线，若需要表达导航性则在被导航端使用开放箭头。表\ref{tab:uml-relations}把线型、端点符号和适用语义放在一起，绘图时先确定关系语义，再选择符号，不要把所有连线都画成普通实线箭头。

\begin{table}[htbp]
\centering
\caption{UML类图关系符号与语义}
\label{tab:uml-relations}
\begin{tabular}{p{2.4cm}p{4.2cm}p{7.0cm}}
\toprule
关系 & 图形符号 & 语义与判读要点 \\
\midrule
泛化 & 实线 + 空心三角 & 子类继承父类，三角指向更一般的类 \\
实现 & 虚线 + 空心三角 & 类实现接口，三角指向接口 \\
关联 & 实线，可带开放箭头 & 对象之间存在结构联系，箭头只表示导航方向 \\
依赖 & 虚线 + 开放箭头 & 使用方在某次操作中依赖被使用方 \\
聚合 & 实线 + 空心菱形 & 整体与部分可独立存在，菱形在整体端 \\
组合 & 实线 + 实心菱形 & 部分的生命周期受整体约束，菱形在整体端 \\
\bottomrule
\end{tabular}
\end{table}

图\ref{fig:uml-class}把“测点产生监测记录、监测记录可触发告警事件”的领域关系落实为类、属性、操作和多重性。读图时先检查职责，再检查关联方向和数量约束。

\begin{figure}[H]
\centering
\begin{tikzpicture}[>=Stealth, node distance=1.0cm and 1.2cm,
  class/.style={draw=Primary, fill=PrimaryLight, text width=3.6cm, align=left, font=\small}]
  \node[class] (station) {\textbf{MonitoringPoint}\\\rule{\linewidth}{0.4pt}\\pointId: String\\type: PointType\\\rule{\linewidth}{0.4pt}\\updateStatus()};
  \node[class, right=of station] (record) {\textbf{MonitoringRecord}\\\rule{\linewidth}{0.4pt}\\timestamp: Instant\\value: Decimal\\quality: QualityCode\\\rule{\linewidth}{0.4pt}\\validate()};
  \node[class, below=of record] (alarm) {\textbf{AlarmEvent}\\\rule{\linewidth}{0.4pt}\\level: AlarmLevel\\status: AlarmStatus\\\rule{\linewidth}{0.4pt}\\acknowledge()};
  \draw[-{Stealth[open]}, thick] (station) -- node[near start, above, font=\scriptsize] {1} node[near end, above, font=\scriptsize] {0..*} (record);
  \draw[-{Stealth[open]}, thick] (record) -- node[near start, right, font=\scriptsize] {1..*} node[near end, left, font=\scriptsize] {0..*} (alarm);
\end{tikzpicture}
\caption{监测点、监测记录与告警事件类图示例}
\label{fig:uml-class}
\end{figure}

\paragraph{（2）活动图（Activity Diagram）}活动图用于描述系统或业务流程中的工作流或操作步骤。它展示一系列动作、决策和控制流之间的关系。活动图用于分析业务逻辑、优化流程以及处理并发行为。

\paragraph{活动（Activity）}表示系统中的操作或步骤，用圆角矩形表示（注意与用例图中的椭圆区分）。

\paragraph{决策节点（Decision Node）}表示条件判断或选择，通常用菱形表示。

\paragraph{并行处理（Fork/Join）}活动图可以表示多个操作并行进行，分叉节点（Fork）表示并行开始，汇合节点（Join）表示并行结束。

\paragraph{（3）状态图（State Diagram）}状态图表示对象或系统在其生命周期中的各种状态及状态之间的转移条件。状态图通常用于描述具有复杂行为的对象或系统，特别是状态驱动的系统。

\paragraph{状态（State）}表示系统或对象在某一时刻的某一状态，通常用圆角矩形表示。

\paragraph{事件（Event）}表示触发状态转换的事件，通常用箭头表示。

\paragraph{状态转换（Transition）}表示从一个状态到另一个状态的转换，通常用带箭头的实线表示。

状态图对于系统中如告警处置、巡检工单流转等具有多个阶段的过程非常有用。

告警事件是案例水库平台中适合绘制状态图的对象：它从规则服务创建开始，经过通知、确认和处置，最终闭环或被专业分析员判为误报。图\ref{fig:alarm-event-state}把正常路径和超时升级放在同一张图中，读者应特别注意“已通知”到“处置中”的超时升级转移并不会跳过审批责任，升级事件仍要留下原事件号和升级原因。

\begin{figure}[H]
\centering
\begin{tikzpicture}[node distance=7mm and 9mm,
  state/.style={draw=Primary,rounded corners,fill=PrimaryLight,minimum width=2.2cm,
    minimum height=.68cm,align=center,font=\small},
  trans/.style={->,draw=Primary,thick}]
\node[state] (new) {新建};
\node[state,right=of new] (notified) {已通知};
\node[state,right=of notified] (acked) {已确认};
\node[state,below=of acked] (handling) {处置中};
\node[state,left=of handling] (closed) {已闭环};
\node[state,right=of handling] (false) {已误报};
\draw[trans] (new)--node[above,font=\scriptsize]{发送通知}(notified);
\draw[trans] (notified)--node[above,font=\scriptsize]{确认}(acked);
\draw[trans] (acked)--node[right,font=\scriptsize]{生成工单}(handling);
\draw[trans] (handling)--node[below,font=\scriptsize]{回执与复核}(closed);
\draw[trans] (acked)--node[above right,font=\scriptsize]{判定误报}(false);
\draw[trans,dashed] (notified) to[bend left=25] node[above,font=\scriptsize]{超时升级}(handling);
\draw[trans,dashed] (handling) to[bend right=25] node[below,font=\scriptsize]{复核为误报}(false);
\end{tikzpicture}
\caption{告警事件生命周期状态图}
\label{fig:alarm-event-state}
\end{figure}

状态转换必须由事件、权限和证据共同决定。例如“已通知”不能仅因页面刷新就变成“已确认”，而要有值班员身份、确认时间和事件版本；“已闭环”需要处置回执与专业复核，“已误报”要保留误报原因并进入规则改进清单。超时升级可以改变责任岗位和通知渠道，但不能覆盖原状态历史。这样的状态模型能支撑幂等处理、重复确认和审计回放。

活动图则强调“谁在什么泳道执行哪一步”。图\ref{fig:alarm-activity}将预警处置拆成平台和岗位两个泳道，展示事件创建、通知、确认、审批和回执的顺序；与状态图相比，它关注动作流和责任分工，而不是对象当前状态。

\begin{figure}[H]
\centering
\begin{tikzpicture}[node distance=6mm and 8mm,
  act/.style={draw=Primary,rounded corners,fill=PrimaryLight,minimum width=2.55cm,
    minimum height=.62cm,align=center,font=\small},
  lane/.style={draw=TextGray,dashed,minimum width=5.9cm,minimum height=4.25cm},
  flow/.style={->,draw=Primary,thick}]
\node[lane] (platformLane) at (-3.1,0) {};
\node[lane] (roleLane) at (3.1,0) {};
\node[font=\small\bfseries] at (-3.1,1.85) {平台};
\node[font=\small\bfseries] at (3.1,1.85) {值班员/审批人};
\node[act] (create) at (-3.1,1.0) {创建事件};
\node[act] (notify) at (-3.1,0.0) {发送通知};
\node[act] (ack) at (3.1,0.0) {确认并研判};
\node[act] (approve) at (3.1,-1.0) {审批处置};
\node[act] (receipt) at (-3.1,-1.0) {记录回执};
\draw[flow] (create)--(notify);
\draw[flow] (notify)--(ack);
\draw[flow] (ack)--(approve);
\draw[flow] (approve)--(receipt);
\end{tikzpicture}
\caption{预警处置泳道活动图}
\label{fig:alarm-activity}
\end{figure}

\paragraph{（4）时序图（Sequence Diagram）}时序图按时间顺序描述参与者与构件之间的消息交互，适合核对跨模块责任、同步与异步边界以及异常分支。图\ref{fig:uml-sequence}展示“告警触发—推送—确认”的主路径：分析服务生成事件，消息服务推送通知，值班员确认后由告警服务持久化状态并返回结果。实际设计还应补充推送失败、重复确认和超时升级分支。

\begin{figure}[H]
\centering
\begin{tikzpicture}[>=Stealth, every node/.style={font=\small}]
  \foreach \x/\name in {0/分析服务,3/告警服务,6/消息服务,9/值班终端}{
    \node[draw=Primary, fill=PrimaryLight, rounded corners=2pt, minimum width=2.0cm] (n\x) at (\x,0) {\name};
    \draw[dashed, draw=TextGray] (\x,-0.4) -- (\x,-5.0);
  }
  \draw[draw=Primary, fill=PrimaryLight] (-0.08,-1.0) rectangle (0.08,-1.8);
  \draw[draw=Primary, fill=PrimaryLight] (2.92,-1.9) rectangle (3.08,-4.0);
  \draw[draw=Primary, fill=PrimaryLight] (5.92,-2.8) rectangle (6.08,-3.4);
  \draw[draw=Primary, fill=PrimaryLight] (8.92,-2.8) rectangle (9.08,-4.0);
  \draw[->, thick] (0,-1.0) -- node[above] {创建告警事件} (3,-1.0);
  \draw[-{Stealth[open]}, thick] (3,-1.9) -- node[above, fill=white, inner sep=1pt] {异步推送} (6,-1.9);
  \draw[-{Stealth[open]}, thick] (6,-2.8) -- node[above, fill=white, inner sep=1pt] {告警通知} (9,-2.8);
  \draw[->, thick] (9,-3.7) -- node[above] {确认（事件号、版本）} (3,-3.7);
  \draw[dashed,-{Stealth[open]}, thick] (3,-4.6) -- node[above, fill=white, inner sep=1pt] {确认结果} (9,-4.6);
\end{tikzpicture}
\caption{告警触发、推送与确认时序图（实心箭头为同步消息，开放箭头为异步或返回消息）}
\label{fig:uml-sequence}
\end{figure}

在案例水库中，几种视图不是彼此独立的插图，而是一组可以相互校验的设计证据。建模评审应先从用例图开始：值班员确认预警、专业分析员维护规则分别对应哪些业务目标，系统边界外的角色是否都能找到责任人。随后把每个重要用例落实到类图中的职责，例如告警服务负责生成和更新事件，监测记录负责保存测值与质量码，监测点只维护自身标识和采集状态。若一个类同时承担通知、审批、历史查询等互不相干的职责，就应拆分服务或引入协作者，而不应靠增加更多方法掩盖耦合。

类图中的多重性还要和数据约束、接口返回值保持一致。一个监测点可以产生零到多条记录，记录可以不触发告警，也可以在不同规则下触发多个事件；因此图中的数量不能被实现成一个必填的单值字段。对告警事件使用 $1..*$ 的一端，表示每个事件至少关联一条触发依据，审计时可以沿着事件号追溯到原始记录。若业务规则允许人工创建没有测值依据的事件，应在需求和数据库约束中明确例外，而不是悄悄把图上的多重性改成 $0..*$。

状态图、活动图和时序图则从三个角度检查同一条处置链。状态图回答“事件此刻处于什么阶段、哪些转移合法”，活动图回答“哪一个岗位执行下一步动作”，时序图回答“跨服务消息以何种顺序到达”。例如确认操作在状态图中只能把“已通知”变为“已确认”，在活动图中由值班员泳道执行，在时序图中则必须携带事件号和版本号写回告警服务。三张图若出现不一致，优先回到业务规则核对；不要为了让箭头看起来连贯而删除超时、拒绝或重复消息等异常分支。

图形符号也应服务于审查，而不是装饰。同步调用需要等待返回结果，适合用于创建事件和保存确认；异步消息只保证进入消息通道，适合用于通知推送，后续应由重试、幂等键和死信处理兜底。开放箭头、虚线返回线和激活条可以帮助读者判断等待关系，但它们不能替代接口契约。实现时仍需写清超时、重试次数、事件版本和操作者身份，并把这些字段作为日志或审计记录的一部分保存。

提交 UML 文档前可按以下顺序自检：第一，所有角色、用例和岗位名称是否采用全书统一术语；第二，每一条关联是否说明方向、数量和生命周期归属；第三，正常、超时、拒绝、重复提交和误报路径是否至少在一种行为视图中出现；第四，图中的类名、方法名与代码接口是否一致；第五，状态转移是否有触发事件、权限条件和可追溯证据。通过这五项检查，图纸才能成为开发、测试和运维共同使用的契约，而不是只在评审会上展示的静态图片。

\paragraph{部署模型}

部署模型用于描述系统硬件、软件构件的物理部署以及它们之间的网络配置。它确保了系统架构能够有效实施，并且能够支持系统的性能、扩展性、可靠性等非功能性需求。部署模型常用于分析系统的分布式架构和硬件资源配置，确保在实际环境中能够满足系统的需求。

\paragraph{节点（Node）}节点代表物理硬件或执行环境，它是部署图的基本元素。节点可以是服务器、客户端、数据库或其他硬件设备。在部署模型中，节点通常用矩形表示，包含该节点上的组件或软件实例。

\paragraph{组件（Component）}在部署模型中，组件是运行在某个节点上的软件构件，通常表示为矩形框，框内标注组件的名称。

\paragraph{通信关系（Communication）}节点之间的通信关系通常用带箭头的实线表示，箭头指向消息传递的方向。通信关系描述了不同节点之间如何交换数据，如何进行同步或异步通信。

\paragraph{部署关系}部署关系用来表示节点和组件之间的关系，指示哪些组件部署在具体的节点上。

通过部署模型，开发人员可以清晰地了解系统中各个组件在物理设备上的部署情况，并根据需要进行硬件资源的规划和优化。部署模型对于分布式系统尤为重要，因为它能帮助理解如何在多台机器或多个数据中心之间分配负载、确保冗余和容错。

\paragraph{构件图}

构件图用于展示软件系统中的构件（模块）及其依赖关系。构件图可以帮助开发人员理解系统如何分解为不同的功能单元，支持模块化设计和集成。通过构件图，系统的组成部分和它们之间的依赖关系可以清晰地呈现出来。

\paragraph{构件（Component）}构件是系统中的一个独立模块，通常用于表示某个功能模块或服务单元。在构件图中，构件用矩形框表示，框内标明构件的名称。

\paragraph{接口（Interface）}接口定义了构件提供的功能或服务，或者构件所需要的功能。在构件图中，接口通常用小圆圈表示，并与构件通过依赖关系或实现关系连接。

\paragraph{依赖关系（Dependency）}依赖关系表示一个构件依赖于另一个构件的接口或服务，通常通过带箭头的虚线表示。

\paragraph{提供与要求接口}构件可以提供某些服务（通过提供接口），也可以依赖于其他构件提供的服务（通过要求接口）。

\subsubsection{架构决策的核心要素}
架构决策需综合考虑以下核心因素，确保设计的合理性与可行性：

\paragraph{功能需求}

明确系统必须实现的核心功能是架构决策的首要步骤，它直接决定了系统的基本架构和模块划分。以水利工程安全监测平台为例，其核心功能包括监测数据接入、异常识别和预警处置等。平台需要持续接入渗压计、位移计、GNSS 监测站、雨量站、库水位计以及视频监控等多类设备数据，并将其传输到中心平台；当监测指标超过阈值或趋势分析识别到异常时，平台要能够迅速触发预警、推送告警信息并支撑值班员开展研判和处置，这就要求系统具备实时处理能力、可靠通信机制和稳定的告警链路。

\paragraph{非功能需求}

非功能需求决定系统在真实负载和故障条件下是否仍可用。性能方面，响应时间需按具体操作和负载定义，要求系统在硬件选型、软件算法优化以及网络架构设计等方面提高数据处理速度和传输效率。可靠性上，系统需支持 24/7 运行，需要高可用性架构，包括冗余设计、故障转移机制和自我检测修复功能。安全性方面，数据传输需加密，需采用加密算法对传输数据进行处理，设置访问控制和身份验证机制，防止未经授权的用户获取或篡改数据。

\paragraph{技术选择}

技术选择是架构决策中的重要环节，影响系统的开发效率、性能表现和可维护性。本书案例锁定 Java 17、Spring Boot 3.2+、Spring Security 6、PostgreSQL + PostGIS + TimescaleDB、Redis、Kafka 以及 Vue 3.4+；若项目需要替换其中任一组件，应通过 ADR 记录兼容性、迁移成本、性能和安全评测证据。在选择技术栈时，需要平衡开发效率与系统性能，既要确保开发过程顺利，又要保证系统在上线后的稳定运行。

四、风险与约束

技术风险和业务约束必须在架构决策中形成可验证的缓解措施。第三方模块兼容性可能导致集成异常或性能下降，因此应在项目前期建立兼容性矩阵和验证环境，并跟踪安全与版本更新；预算和时间限制则要求按业务价值与风险排列交付顺序。风险清单还应记录责任人、触发条件、应对动作和复核日期，使风险控制能够被检查，而不只停留在原则表述。

五、扩展性与复用性

扩展性与复用性决定了架构在未来的适应空间。案例水库平台若要新增一类光纤渗压传感器，理想情况下只需新写一个协议适配器并在设备台账登记，而不改动质量检查、预警评估和展示层——能否做到，取决于采集接入是否被隔离在稳定接口之后。复用性同理：认证与审计模块若不掺入水库特有逻辑，即可原样移植到灌区或河道项目。架构评审时应各举一个“下一年最可能发生的扩展”来检验决策，而不是抽象地宣称系统“可扩展、可复用”。

\subsection{架构决策记录（ADR）}

架构决策记录（Architecture Decision Record，ADR）是把一次重要技术取舍固定为可追踪证据的短文档。它不替代需求规格，而是说明为什么在当时的约束下选择某个方案，以及未来什么条件变化时需要重新评估。每条 ADR 至少包含：编号、标题、状态、上下文、决策、备选方案、后果、关联需求、日期和责任人。状态可取“提议、已接受、已取代、已废弃”，取代关系必须链接到新记录。

以告警链路为例，ADR-003 的上下文是汛期数据入站峰值升高且通知服务可能短时不可用；决策是在分析服务与通知服务之间使用持久化消息队列，事件写入成功后再异步推送。备选方案是同步调用（实现简单但下游故障会阻塞生成）和仅内存队列（吞吐高但重启会丢事件）。接受该决策的后果是增加消息代理、幂等键、积压监控和回放运维成本，同时获得可重试和故障恢复能力。关联需求为 REQ-MON-03（告警生成 P95 不超过 5 秒、推送送达 P95 不超过 10 秒）与 REQ-MON-04（告警事件可审计）；责任人和日期应由项目组填写，评审证据放在记录末尾。

ADR 的评审顺序是先确认上下文和需求，再比较备选方案的质量属性影响，最后列出验证任务和回滚条件。若后续将消息队列替换为其他实现，不应直接修改旧记录，而应新建 ADR，标记旧记录“已取代”，并在迁移测试、数据回放和运维手册中引用新编号。这样，架构演化与代码变更、部署变更保持同一条审计链。

反面样本同样有教学价值。一份写坏的 ADR 通常长这样：“决策：使用Kafka。理由：Kafka性能好、业界主流。”它缺了上下文（为什么此处需要队列）、缺了备选（同步调用与内存队列为何不行）、缺了后果（多出哪些运维成本）、缺了重估条件（什么变化时该重看）。这样的记录在半年后毫无价值——没人能判断当时的前提是否仍然成立。判断一份 ADR 合格与否有个简单测试：把结论遮住，只看上下文与备选方案分析，读者应能大致推出同样的结论；推不出来，说明记录的是立场而不是决策。

\section{架构设计原则}

软件架构设计需要遵循一系列基本原则，这些原则指导我们构建高质量、可维护、可扩展的软件系统。主要原则包括模块化与关注点分离、抽象与信息隐藏、功能独立性与复用性等；模式文献则为反复出现的设计问题提供可复用词汇\cite{martin2017,gamma1994}。下面将逐一介绍这些核心原则。

\subsection{模块化与关注点分离原则}

模块化与关注点分离是软件架构设计的核心原则，旨在通过将系统分解为独立、可管理的模块，降低复杂性并提升可维护性。本小节讨论模块划分的内聚性与耦合性，以及如何通过这些原则优化系统设计。

\subsubsection{模块化的定义与目标}
模块化（Modularity）是将软件系统分解为多个独立模块的过程，每个模块负责特定的功能或子功能。其核心目标是：

\paragraph{降低复杂性}将复杂系统拆分为较小的单元，便于开发、测试和维护。例如，一个智慧水利平台可以被拆分成监测接入、数据治理、告警处置、巡检管理等多个模块，每个模块都可以独立进行开发和测试，从而降低了整个系统的复杂性。

\paragraph{提高复用性}独立模块可在不同场景中复用，减少重复开发。例如，一个通用的用户认证模块可以被多个不同的应用程序所使用，无需每次重新开发相同的认证功能。

\paragraph{支持并行开发}团队成员可独立开发不同模块，提升协作效率。例如，在一个软件开发项目中，不同的开发人员各自负责不同的模块开发，如前端界面、后端逻辑、数据库交互等，从而提高整个团队的开发效率。

示例：在水利工程安全监测平台中，可将“感知接入”“预警分析”“可视化交互”划分为独立模块，每个模块负责特定功能。感知接入模块负责与各类监测设备通信并汇聚原始数据；预警分析模块根据监测数据、阈值规则和趋势模型判断是否触发预警；可视化交互模块则为用户提供监测总览、趋势曲线和告警处置界面。

\subsubsection{关注点分离（Separation of Concerns）}
关注点分离（SoC）是一种设计哲学，强调将系统中不同的“关注点”（如业务逻辑、用户界面、数据存储）分离到独立模块中。其核心思想是：

\paragraph{单一职责原则}每个模块应专注于解决一个特定的问题领域。例如，在一个内容管理系统中，用户管理模块只负责处理用户相关的操作，如注册、登录、权限管理等；文章管理模块只负责处理文章的创建、编辑、发布等操作。

\paragraph{解耦依赖}减少模块间的相互影响，降低变更成本。例如，如果用户管理模块和文章管理模块是分离的，那么对用户管理模块的修改不会影响到文章管理模块，反之亦然，从而降低了系统变更带来的风险和成本。

\paragraph{示例}水利工程安全监测平台中的“用户认证”模块应独立于“预警判定”模块，确保认证机制的变更不影响监测分析功能。用户认证模块负责验证值班员、专业分析员、审批人和运维员的身份信息；预警判定模块则根据监测数据和预设规则来判断是否触发预警，两者相互独立，互不干扰。

\subsubsection{内聚性（Cohesion）}
内聚性衡量模块内部各元素（如功能、数据）的关联程度。高内聚的模块功能单一，易于理解和维护。以下是常见的内聚类型：

\paragraph{功能内聚（Functional Cohesion）}模块内所有元素共同完成单一、明确的功能。例如，在一个图像处理软件中，图像旋转模块只负责对图像进行旋转操作，不涉及其他如图像缩放、滤镜应用等功能，这样的模块具有高度的功能内聚性。

\paragraph{顺序内聚（Sequential Cohesion）}模块内元素按顺序执行，前一操作的输出是后一操作的输入。例如，在一个数据处理系统中，数据清洗模块先对原始数据进行清洗，去除噪声和错误数据，然后将清洗后的数据传递给数据分析模块进行进一步的分析，这两个模块在功能上具有顺序关系，形成了顺序内聚。

\paragraph{通信内聚（Communicational Cohesion）}模块内元素操作相同的数据或资源。例如，告警管理模块中的各个子模块都操作告警数据，如告警生成子模块负责创建告警事件，告警更新子模块负责更新处置状态，告警查询子模块负责查询告警信息，它们都围绕告警数据进行操作，具有通信内聚性。

\paragraph{偶然内聚（Coincidental Cohesion）}模块内元素无逻辑关联，仅因偶然原因组合。例如，将“用户认证”和“报表生成”合并为一个模块，这两个功能本身没有直接的逻辑关联，只是偶然地被放在一起，导致功能混乱，不利于模块的理解和维护。

实践建议：优先追求功能内聚，避免偶然内聚。在模块设计时，应尽量使模块内的元素共同完成一个单一的功能，这样可以提高模块的内聚性，使模块更易于理解和维护，也能提高系统的可扩展性和可修改性。

\subsubsection{耦合性（Coupling）}
耦合性衡量模块之间的依赖程度。低耦合的系统更灵活，易于修改和扩展。以下是常见的耦合类型：

\paragraph{数据耦合（Data Coupling）}模块间通过参数传递数据。例如，在一个文本编辑器中，文本输入模块将用户输入的文本数据通过参数传递给文本显示模块，文本显示模块根据接收到的文本数据进行显示操作，两者之间通过数据参数进行耦合，耦合度较低。

\paragraph{控制耦合（Control Coupling）}模块间通过控制参数传递指令（如标志位）。例如，在一个图形绘制软件中，绘图模块根据主模块传递的控制参数（如绘制直线、绘制圆形等标志位）来决定绘制何种图形，主模块对绘图模块的控制通过控制参数实现，形成控制耦合，这种耦合度比数据耦合稍高。

\paragraph{公共耦合（Common Coupling）}模块共享全局数据或资源。例如，在一个多线程应用程序中，多个线程模块共享同一个全局配置文件，当一个线程模块修改了全局配置文件中的某个参数，其他线程模块也会受到影响，这种共享全局数据的方式导致模块间的耦合度较高，容易引发不可预知的问题。

\paragraph{内容耦合（Content Coupling）}模块直接修改或依赖另一模块的内部数据。

例如，一个模块直接访问另一个模块的私有成员变量并对其进行修改，这种耦合方式模块间的依赖关系非常紧密，一个模块的修改很可能导致另一个模块的错误，耦合度最高，应尽量避免。

实践建议：优先使用数据耦合，避免控制耦合和内容耦合。在设计模块间的交互时，应尽量通过参数传递数据，减少模块间的控制依赖和直接的数据修改，以降低耦合度，提高系统的灵活性和可维护性。

\subsubsection{平衡内聚与耦合}
理想的模块设计应追求高内聚、低耦合。

高内聚确保模块功能单一，易于维护。当模块内的元素紧密协作完成一个特定的功能时，模块的内聚性高，开发人员在理解和维护该模块时会更加容易，因为只需要关注模块内部的有限功能范围。

低耦合减少模块间依赖，降低变更风险。如果模块间的耦合度低，那么一个模块的修改不太可能影响到其他模块，系统的变更风险降低，扩展和修改也更加容易。

\paragraph{案例分析：水利工程安全监测平台}

\paragraph{高内聚}将“异常识别算法”封装为独立模块，仅负责逻辑判断，不涉及数据存储或用户交互。异常识别模块专注于根据位移、渗流、雨量和库水位等监测数据进行分析，判断工程状态是否异常，而不涉及如何存储数据或如何展示页面，这样的设计使算法模块具有高内聚性，便于开发、测试和后续优化。

\paragraph{低耦合}“可视化界面”模块通过 API 与“监测接入模块”交互，无需了解其内部实现。可视化界面模块通过调用标准化数据服务接口获取监测数据，并将其展示为曲线、报表和三维场景，而不需要关心底层采集协议和设备驱动细节，这种通过 API 进行交互的方式降低了界面模块与接入模块之间的耦合度，使得两者可以独立开发和维护。

\subsubsection{实践步骤}
\paragraph{识别关注点}通过需求分析确定系统的核心功能与非功能需求。在项目启动初期，与用户、业务分析师等利益相关者进行深入沟通，了解系统的业务目标、用户需求以及各种非功能需求（如性能、安全性、可靠性等），从而明确系统中的各个关注点，为后续的模块划分提供依据。

\paragraph{划分模块}根据功能相关性和内聚性原则，将系统分解为独立模块。根据之前识别的关注点，将具有紧密功能相关性的元素划分为同一个模块，并遵循内聚性原则，确保每个模块都具有单一、明确的功能职责，避免模块功能过于分散或交叉。

\paragraph{定义接口}明确模块间的交互方式，使用数据耦合传递参数。为每个模块定义清晰的接口，规定模块间的输入输出参数、调用方式等交互细节，尽量采用数据耦合的方式进行模块间通信，减少模块间的依赖关系，提高模块的独立性和可替换性。

\paragraph{评估与优化}使用工具（如依赖分析工具）检测高耦合模块，并重构优化。在开发过程中，利用依赖分析工具对模块间的依赖关系进行检测，找出存在高耦合度的模块对，分析其耦合原因，并通过重构的方式降低耦合度，如调整模块接口、引入中间层、解耦模块间的直接依赖等，不断优化模块设计，以达到高内聚、低耦合的理想状态。

\subsection{抽象与信息隐藏原则}

抽象与信息隐藏是软件架构设计的核心原则，通过层次化抽象简化复杂系统，并通过接口封装隔离实现细节。本小节讨论其定义、应用及实践方法。

\subsubsection{抽象的定义与作用}
抽象（Abstraction）是对复杂系统的简化表示，通过忽略非本质细节，聚焦于核心特征。其核心作用是：

\paragraph{1.降低认知负担}允许开发者以更高层次的视角理解系统，避免陷入细节。例如，在设计一个图形用户界面时，开发者可以将按钮、文本框等界面元素抽象为通用的控件类，而无需在设计界面布局时考虑每个控件内部的绘图算法和事件处理机制，从而能够更专注于界面的整体结构和用户体验。

\paragraph{2.支持复用}抽象层可独立于具体实现，便于在不同场景中复用。例如，一个抽象的数据库访问层可以被多个不同的业务模块所使用，无论这些业务模块是处理测站台账、监测数据还是告警记录，都可以通过相同的抽象接口进行数据库操作，提高了代码的复用性。

\paragraph{3.促进模块化}抽象为模块划分提供依据，确保模块职责明确。例如，在一个智慧水利平台中，可以将用户认证、监测接入、告警分析、三维展示等抽象为不同的模块，每个模块都有明确的职责和接口，便于团队协作开发和后续的维护扩展。

\paragraph{示例}在水利工程安全监测平台中，“监测设备”可抽象为具有“采集”和“上报”功能的对象，而不关心其具体型号或通信协议。这种抽象使得开发人员可以将位移计、渗压计、雨量计、GNSS 终端等不同设备统一视为标准数据源，简化系统设计与实现，便于后续扩展新型监测设备。

\subsubsection{抽象层次的一种划分示例}\label{sec:hierarchical-abstraction}
层次化抽象将系统分解为多个层次，每个层次处理特定的抽象级别，形成“自顶向下”的结构。常见的层次划分包括：

\paragraph{应用层（Application Layer）}编排用例和事务边界，例如接收查询、组织质量校验、调用领域策略并返回结果。应用层不持有设备协议和核心判定规则，而是负责流程协调。

\paragraph{领域/业务层（Domain/Business Layer）}定义安全监测、质量判定和预警策略等业务规则。渗流异常识别、变形趋势判定和告警等级计算都属于这一层，规则只依赖领域对象与接口，不直接访问数据库或设备。

\paragraph{基础设施层（Infrastructure Layer）}提供存储、通信、消息代理和设备驱动等技术能力。它实现数据持久化、网络传输和协议适配，但不决定业务规则；例如 PostgreSQL/TimescaleDB 适配器、Kafka 发布器和传感器驱动都属于该层。

实践建议：

层次间按“应用层（编排用例）→领域/业务层（规则与策略）→基础设施层（存储、通信、设备驱动）”方向依赖，应用层通过端口访问基础设施能力，不能直接调用数据库驱动；这样可以保持各层独立，降低耦合度。

每层仅关注其抽象级别内的细节：业务层只关心规则，基础设施层只关心技术实现。第1章的“感知—平台—应用—用户”是业务视角；本节是软件实现的抽象层次示例，后续软件分层统一采用图\ref{fig:layered-architecture}的口径。

\subsubsection{信息隐藏（Information Hiding）}
信息隐藏通过封装实现细节，仅暴露必要接口，确保模块内部变化不影响外部调用\cite{parnas1972}。其核心原则是：

\paragraph{封装（Encapsulation）}将数据和操作封装在模块内部，外部通过接口访问。例如，在一个水库安全监测系统中，测站管理模块将测站信息、设备状态、采集策略等数据以及新增测站、启停设备等操作封装在内部，其他模块只能通过测站管理模块提供的接口访问相关能力，而不能直接修改其内部数据结构。

\paragraph{接口隔离（Interface Segregation）}为不同客户端提供专用接口，避免不必要的依赖。例如，对于智慧水利平台，可以为值班员提供告警查看与确认接口，为运维员提供设备维护与通信配置接口。这样可以减少不同角色对无关功能的依赖，提高系统的灵活性和可维护性。

\paragraph{示例}水利工程安全监测平台的“监测接入模块”封装了设备通信细节，仅向外部暴露 readMonitoringData() 接口，其他模块无需了解具体实现。模块内部支持 Modbus、MQTT 和厂商私有协议，但这些细节都被隐藏在模块内部，外部模块只需调用统一接口即可获取监测数据。更换采集设备或接入协议时，只要接口不变，其他模块无需修改。

\subsubsection{接口封装的实践方法}
\paragraph{定义清晰的接口}

使用 UML 类图或接口定义语言（如 IDL）明确接口的方法、参数和返回值。例如，在设计一个网络通信模块时，可以使用UML类图定义一个NetworkService接口，其中包含connect()、sendData()、receiveData()等方法，每个方法的参数和返回值都明确指定，如connect()方法接受服务器地址和端口号作为参数，返回一个布尔值表示连接是否成功。

下面用与本章构件级设计一致的接口名定义统一监测数据读取能力：

\begin{lstlisting}[label={lst:ch03-java-interface},language=Java,caption={监测数据提供接口}]
public interface MonitoringDataProvider {
    MonitoringData readMonitoringData();
}
\end{lstlisting}

接口契约还应把输入、输出、错误和版本写成可编译的类型。下面的 DTO 只承载跨层传输字段，业务对象仍由领域层负责不变量校验；代码清单\ref{lst:ch03-monitoring-dto}展示了时间、质量码和规则版本的最小契约。
\begin{lstlisting}[label={lst:ch03-monitoring-dto},language=Java,caption={监测数据传输对象}]
import java.math.BigDecimal;
import java.time.Instant;
public record MonitoringDataDto(String pointId,
        Instant occurredAt, BigDecimal value,
        String qualityCode, String ruleVersion) {
}
\end{lstlisting}

应用层通过端口依赖基础设施，而不是直接依赖数据库驱动。代码清单\ref{lst:ch03-record-port}把保存记录的能力抽象为接口，便于用内存替身进行单元测试。
\begin{lstlisting}[label={lst:ch03-record-port},language=Java,caption={监测记录持久化端口}]
import java.util.Optional;
public interface MonitoringRecordPort {
    MonitoringRecord save(MonitoringRecord record);
    Optional<MonitoringRecord> findLatest(String pointId);
}
\end{lstlisting}

用例服务负责事务编排和权限边界，领域策略则负责业务判断。代码清单\ref{lst:ch03-usecase-service}示范了接收测值、保存记录和发布分析事件的顺序。
\begin{lstlisting}[label={lst:ch03-usecase-service},language=Java,caption={测值接收用例服务}]
import java.time.Instant;
import java.math.BigDecimal;
public final class ReceiveMonitoringUseCase {
    private final MonitoringRecordPort records;
    private final AnomalyEventPublisher publisher;
    public ReceiveMonitoringUseCase(MonitoringRecordPort records,
            AnomalyEventPublisher publisher) {
        this.records = records; this.publisher = publisher;
    }
    public void handle(MonitoringData data) {
        MonitoringRecord saved = records.save(MonitoringRecord.from(data));
        publisher.publish(AnomalyEvent.from(saved));
    }
}
\end{lstlisting}

控制器只负责协议转换和参数校验，不把预警规则写在 Web 层。代码清单\ref{lst:ch03-controller}以 Spring MVC 的构造器注入保持依赖显式。
\begin{lstlisting}[label={lst:ch03-controller},language=Java,caption={监测数据查询控制器}]
import org.springframework.web.bind.annotation.GetMapping;
import org.springframework.web.bind.annotation.PathVariable;
import org.springframework.web.bind.annotation.RequestMapping;
import org.springframework.web.bind.annotation.RestController;
@RestController
@RequestMapping("/api/monitoring")
public class MonitoringController {
    private final MonitoringQuery query;
    public MonitoringController(MonitoringQuery query) { this.query = query; }
    @GetMapping("/{pointId}/latest")
    public MonitoringDataDto latest(@PathVariable String pointId) {
        return query.readLatest(pointId);
    }
}
\end{lstlisting}

事件契约需要稳定的幂等标识和版本字段。代码清单\ref{lst:ch03-anomaly-event}用不可变记录表达发布内容，消费者据此实现重复消息检测和规则换版兼容。
\begin{lstlisting}[label={lst:ch03-anomaly-event},language=Java,caption={异常分析事件契约}]
import java.math.BigDecimal;
import java.time.Instant;
public record AnomalyEvent(String eventId, String pointId,
        Instant occurredAt, String qualityCode,
        String ruleVersion, BigDecimal score) {
}
\end{lstlisting}

\paragraph{隐藏实现细节}

将内部数据和实现逻辑设为私有，仅通过接口提供服务。例如，在监测接入模块的实现中，将设备驱动程序、通信协议处理等内部数据和逻辑设为私有成员，外部模块无法直接访问，只能通过\texttt{MonitoringDataProvider}接口中的方法进行操作。

\paragraph{示例}监测接入模块内部可使用Modbus、MQTT或厂商协议通信，但外部只通过\texttt{MonitoringDataProvider}调用，无需了解数据帧和重连等实现细节。

\paragraph{版本控制}

通过版本化接口（如 API 版本号）管理不兼容变更，确保迁移期的新旧调用方能够共存；兼容性变更优先保持原接口名称和契约稳定。

\subsubsection{案例分析：水利工程安全监测平台的抽象与封装}
\paragraph{层次化抽象}

(1)业务层：定义“预警评估策略”抽象类，包含 evaluateRisk() 方法。该抽象类规定了风险评估的基本框架和接口，不同的具体策略类可以继承并实现该抽象类，例如渗流异常识别策略、变形趋势判定策略等。

(2)应用层：编排具体用例，调用业务层策略接口。应用层的预警分析模块根据业务层定义的接口获取经过校验的监测数据，并通过 evaluateRisk() 方法请求渗流异常识别策略判断是否触发工程预警；策略本身仍归属于业务层。

\paragraph{接口封装}

定义\texttt{MonitoringAcquisitionService}接口，暴露\texttt{pauseAcquisition(pointId, reason, until)}和\texttt{resumeAcquisition(pointId)}方法，隐藏设备协议、采样调度和断线重传逻辑。暂停只用于设备检修等受控场景，必须校验授权、记录原因和到期时间，并自动恢复单个测点采集，不中断整站运行。

抽象与信息隐藏沿用本节前述“识别关注点—划分模块—定义接口—评估优化”步骤，但评审重点转为两项：接口是否泄露设备协议、存储结构等实现细节，以及替换实现后调用方测试是否仍能通过。这样避免重复一套流程，同时把原则落实为可检验条件。

\subsection{功能独立性与复用性原则}

功能独立性由内聚与耦合两把尺子度量。与其背诵七级分类的定义，不如把每一级对应到监测平台里的一个具体形态，表\ref{tab:ch03-cohesion-coupling}给出这份对照：左半部分从偶然内聚到功能内聚由差到好，右半部分从内容耦合到数据耦合由差到好。评审代码时拿着这张表找位置，比抽象讨论“高内聚低耦合”有效得多。

\begin{table}[htbp]
\centering
\caption{内聚与耦合分级在监测平台中的对照}
\label{tab:ch03-cohesion-coupling}
\begin{tabular}{p{2.2cm}p{4.6cm}p{2.2cm}p{4.2cm}}
\toprule
内聚等级 & 平台中的形态（差→好） & 耦合等级 & 平台中的形态（差→好） \\
\midrule
偶然内聚 & “工具类”里混着日期格式化与阈值计算 & 内容耦合 & 预警模块直接改写接入模块的内部变量 \\
逻辑内聚 & 一个“处理数据”函数用标志位分流五种业务 & 公共耦合 & 各模块共享一个可写的全局配置对象 \\
时间内聚 & “系统启动时要做的事”堆在一个初始化类 & 控制耦合 & 调用方传标志位决定被调方走哪条分支 \\
过程内聚 & 按操作顺序而非业务归属组织函数 & 印记耦合 & 只需测点编号却传入整个测点对象 \\
通信内聚 & 操作同一数据集的函数聚在一起 & 数据耦合 & 接口只传必需的值参数 \\
功能内聚 & 渗流异常识别模块只做异常识别 & —— & —— \\
\bottomrule
\end{tabular}
\end{table}

功能独立性与复用性是软件架构设计的核心目标，旨在通过模块化设计降低系统复杂度，提升开发效率与维护性。本小节讨论基于构件的设计原则及其在实践中的应用。

\subsubsection{基于构件的设计原则}
基于构件的设计（Component-Based Design, CBD）通过以下原则实现功能独立性与复用性：

\paragraph{1.领域工程（Domain Engineering）}

领域工程是基于构件的设计的基础，它涉及到对特定领域的深入理解和分析，以提取出可复用的构件。这些构件封装了领域的共性需求和专业知识，可以在不同的项目中被反复使用，从而提高开发效率和系统质量。

\paragraph{实践步骤}

识别领域内的通用功能（如身份验证、日志记录）：在进行领域分析时，需要对领域内的各种应用进行调研和分析，找出其中的通用功能模块。例如，在企业级应用中，身份验证和日志记录几乎是每个系统都必需的功能。这些功能在不同的应用中具有相似的流程和逻辑，可被识别为通用功能模块。

设计通用构件，封装领域知识：一旦识别出通用功能，下一步就是设计可复用的构件来实现这些功能，并将相关的领域知识和业务规则封装在构件内部。例如，设计一个通用的身份验证构件，它包含了验证用户名和密码的逻辑，也封装了与用户管理、权限控制等相关的领域知识。这样，在不同的项目中，只需简单地配置和调用该构件，即可实现符合项目需求的身份验证功能，无需重复开发。

\paragraph{构件合格性检验}

构件合格性检验是确保基于构件的设计成功的重要环节，它通过对构件进行全面的测试和验证，保证构件在功能、性能和质量等方面满足系统的要求。

\paragraph{方法}

\paragraph{单元测试}验证构件内部逻辑。单元测试是对构件最基本的功能单元进行测试，目的是验证构件内部的逻辑是否正确实现。例如，对于一个数据处理构件，单元测试可以检查其对各种输入数据的处理是否准确无误，是否能够正确地进行数据转换、计算等操作。通过单元测试，可以及时发现和修复构件内部的缺陷，确保构件的功能完整性。

\paragraph{集成测试}检查构件与其他模块的兼容性。集成测试是在将构件集成到系统中后进行的测试，主要检查构件与系统中其他模块之间的交互是否正常，是否存在兼容性问题。例如，在一个分布式系统中，集成测试可以验证一个通信构件是否能够与其他模块正确地进行数据传输和交互，重点观察三种失败形态：连接建立不上、报文能发出但字段被截断、以及重传后同一条读数被写入两次。这三种情况在单元测试里都看不出来，只有把构件放回真实系统环境才会暴露。

\paragraph{复用的分析与设计}

复用的分析与设计是基于构件的设计的核心，它关注如何通过合理的接口设计和配置机制，提高构件的通用性和适应性，使其能够在不同的系统和场景中被广泛复用。

\paragraph{接口标准化}设计通用接口（如 REST API），支持跨系统集成。接口一旦标准化，调用方就不必为每个新系统改写一次适配代码。例如，一个监测数据接入构件如果提供了标准化的 REST API 接口，那么无论是监测平台、移动巡检应用还是专题分析系统，都可以方便地调用该构件的数据接口，无需关心其内部实现细节。

\paragraph{配置灵活性}通过参数化配置适应不同场景（如传感器类型可配置）。为了使构件能够在不同的应用场景中灵活使用，需要设计参数化配置机制。例如，在一个传感器数据采集构件中，可以通过配置文件或参数设置来指定传感器的类型、连接方式、采样频率等参数。这样，同一个构件就可以适应不同类型的传感器和不同的工作环境，提高了构件的通用性和复用性。

\subsection{案例分析：水利工程安全监测平台的监测接入模块}

\paragraph{功能独立性}

监测接入模块仅负责数据采集与传输，不参与预警判定逻辑。在水利工程安全监测平台中，该模块专注于采集位移、渗压、雨量、库水位、视频巡检等实时数据，并通过统一协议转换和初步校验后传输给中心平台。这种功能上的独立性使得监测接入模块的职责清晰明确，开发人员在设计和维护该模块时只需关注采集准确性、接入稳定性和协议适配效率，无需考虑预警触发条件判断或用户界面展示等业务逻辑，从而降低系统整体复杂度，提高模块的可维护性和可测试性。

\paragraph{复用性}

通过标准化接口，支持多种监测设备类型（如渗压计、位移计、GNSS、雨量计）。监测接入模块采用标准化接口设计，能够兼容多种不同类型的设备。只要符合模块的接入规范，就可以方便地接入系统。这种设计便于平台扩展和升级，允许项目根据工程实际需求增加新的监测点型，也使得接入模块本身具有较强的通用性和复用性。例如，在大坝安全监测系统、灌区量测系统或河道水情监测系统中，都可以直接复用该模块来采集相应数据，无需重新开发类似能力。

模块可移植至其他水利监测系统。监测接入模块的高复用性还体现在其可移植性上。由于模块在设计时遵循了通用接口标准和通信协议，并且不依赖于特定业务系统的内部实现细节，它可以方便地移植到其他水利信息系统中。例如，在河道水文监测系统中，可以复用该模块采集流量、水位和雨量数据；在泵站运行监测系统中，也可以复用该模块接入振动、电流和视频数据。这种跨系统复用能力提高了软件资源利用率，也促进了不同水利业务系统之间的集成。

\paragraph{与其他模块的交互}

监测接入模块通过标准数据契约把采集结果发送给数据处理服务；数据处理服务完成质量检查和标准化后，分别写入数据存储服务并向告警服务发布可分析事件。告警服务依据已批准规则判断是否触发预警。反向配置由平台通过受控接口下发到监测接入模块，例如调整采样频率或传输参数，并记录操作者、版本和生效时间。这样的职责划分边界明确，也使各服务能够独立测试和升级。

\section{架构类型与常见模式}

\subsection{常见架构类型}

软件架构类型决定了系统的结构、模块交互方式及性能特征。以下是三种最常见的架构类型及其核心要点：

\subsubsection{分层架构（Layered Architecture）}\label{sec:software-layered-architecture}
分层架构将软件实现划分为多个垂直层次，每层专注于特定职责。这里讨论的是软件调用与依赖边界，与第1章“感知—平台—应用—用户”的业务总体架构不是同一划分维度；前者约束代码依赖，后者描述数据上行和控制下行。具体调用方向见图\ref{fig:layered-architecture}，本小节只讨论架构模式及其权衡。企业应用模式文献对分层与数据访问边界有系统讨论\cite{fowler2002}。

\paragraph{核心特点}

\paragraph{职责分离}每层独立演进，降低耦合。分层架构通过将系统划分为不同的层次，使得每一层都能够专注于特定的功能职责。例如，表示层负责与用户进行交互，展示用户界面；业务逻辑层负责处理系统的业务规则和逻辑；数据层则负责数据的存储和访问。由于各层之间的相对独立性，某一层次的变更不会对其他层次产生直接的影响，从而降低了系统整体的耦合度，提高了系统的可维护性和可扩展性。

\paragraph{可扩展性}可单独扩展某一层（如优化数据层存储）。分层架构的另一个重要特点是其良好的可扩展性。由于每一层都具有明确的职责和接口，在需要对系统进行扩展或优化时可以针对特定的层次进行单独的扩展。例如，在数据层，可以通过优化数据库结构、引入缓存机制或采用更高效的存储引擎等方式来提升数据存储和访问的性能，而无需对表示层或业务逻辑层进行大规模的修改。这种单独扩展的能力使得系统能够根据实际需求灵活地进行性能优化或功能扩展，提高了系统的适应性和使用寿命。

\paragraph{典型应用}

Web 应用（如用户界面层、后端逻辑层、数据库层）。在典型的Web应用中，分层架构得到了广泛的应用。用户界面层通常由前端框架构建，负责与用户进行交互，展示页面内容并接收用户输入；后端逻辑层则处理业务逻辑，如用户认证、数据验证、业务流程控制等；数据库层负责数据的持久化存储和查询操作。通过这种分层方式，Web应用的各个部分能够独立开发、测试和部署，提高了开发效率和系统稳定性。

水利工程安全监测平台：将监测接入、预警分析、用户界面分离。在该平台中，分层架构也被有效应用。监测接入层负责与各类监测设备通信，获取原始数据；预警分析层根据监测数据和预设规则进行分析判断，决定是否触发告警；用户界面层则为用户提供监测总览、趋势曲线和告警处置界面。这种层次划分使系统的各个功能模块相互独立，便于开发、测试和后续维护升级。典型的分层架构模式如图\ref{fig:layered-architecture}所示。

\begin{figure}[H]
\centering
\begin{tikzpicture}[>=Stealth, scale=0.9, thick]
  \node[draw=Primary, rectangle, minimum width=5cm, minimum height=0.7cm, fill=PrimaryLight, text=PrimaryDark, rounded corners=2pt, thick] at (0,3) (ui) {\small\bfseries 表示层（用户界面）};
  \node[draw=Primary, rectangle, minimum width=5cm, minimum height=0.7cm, fill=PrimaryLight, text=PrimaryDark, rounded corners=2pt, thick] at (0,2) (biz) {\small\bfseries 业务逻辑层};
  \node[draw=Primary, rectangle, minimum width=5cm, minimum height=0.7cm, fill=PrimaryLight, text=PrimaryDark, rounded corners=2pt, thick] at (0,1) (data) {\small\bfseries 数据访问层};
  \node[draw=Primary, rectangle, minimum width=5cm, minimum height=0.7cm, fill=PrimaryLight, text=PrimaryDark, rounded corners=2pt, thick] at (0,0) (db) {\small\bfseries 数据库层};
  \draw[->,draw=Primary, thick,>=Stealth] (ui.south) -- (biz.north);
  \draw[->,draw=Primary, thick,>=Stealth] (biz.south) -- (data.north);
  \draw[->,draw=Primary, thick,>=Stealth] (data.south) -- (db.north);
\end{tikzpicture}
\caption{软件实现分层的调用依赖方向}
\label{fig:layered-architecture}
\end{figure}

\subsubsection{客户/服务器架构（Client-Server Architecture）}
客户/服务器架构是一种经典的网络架构模式，它将系统分为客户端和服务器端两个部分，通过网络通信进行协作。这种架构模式在许多需要分布式计算和资源共享的场景中得到了广泛的应用。

\paragraph{核心特点}

\paragraph{集中式管理}服务器统一处理数据和业务逻辑。在客户/服务器架构中，服务器端承担了系统的核心数据处理和业务逻辑实现任务。客户端的主要功能是与用户进行交互，收集用户输入并展示服务器端返回的结果。服务器端对所有的数据和业务规则进行集中式的管理和处理，这种集中式管理的方式使得系统的维护和升级更加方便，因为只需要在服务器端进行相应的操作，而无需在每个客户端都进行修改。

\paragraph{资源共享}客户端通过请求-响应模式访问服务器资源。客户端通过向服务器端发送请求，获取所需的资源和数据。服务器端在接收到客户端的请求后，进行相应的处理，并将结果以响应的形式返回给客户端。这种请求-响应模式使得多个客户端能够共享服务器端的资源，避免了资源的重复配置和浪费。例如，在一个省级智慧水利平台中，监测大屏、移动巡检端和管理后台都可以同时访问服务器端的测站和告警数据。

\paragraph{优缺点}

\paragraph{优点}易于维护，资源集中管理。由于服务器端集中处理了系统的数据和业务逻辑，在系统维护和升级时只需对服务器端进行操作，大大简化了维护工作。资源的集中管理提高了资源的利用率，降低了系统的总体成本。服务器端可以提供更强大的计算能力和数据存储能力，支持更多用户的并发访问，提高了系统的性能和可扩展性。

\paragraph{缺点}服务器负载高，单点故障风险。客户/服务器架构的一个主要缺点是服务器端可能面临较高的负载压力，尤其是在用户数量众多或并发请求量较大的情况下。如果服务器端的性能不足或出现故障，将导致整个系统无法正常工作，存在单点故障的风险。为了解决这个问题，通常需要采用负载均衡、集群部署等技术来提高服务器端的性能和可靠性。

\subsubsection{分布式架构（Distributed Architecture）}
分布式架构是一种将系统拆分为多个独立服务的架构模式，这些服务通过网络通信协作完成任务。它的价值在于让高并发、大数据量和差异化的业务模块可以各自独立扩容和发布，代价则是网络调用、部署编排和分布式数据一致性都要由团队自己承担。

\paragraph{核心特点}

\paragraph{可扩展性}通过增加服务实例应对高负载。分布式架构具有很强的可扩展性，可以通过增加服务实例的数量来应对高负载的情况。当系统的访问量或数据量增加时，可以根据实际需求动态地添加新的服务实例，分担现有服务的压力，从而提高系统的整体性能和吞吐量。例如，在汛期高峰时段，可以通过增加监测接入服务、告警分析服务等实例的数量，快速提升系统的处理能力。

容错性：部分服务故障不影响整体系统。分布式架构中的每个服务都是独立运行的，即使其中某个服务出现故障，也不会导致整个系统崩溃。系统可以通过故障转移、自动重启等机制来恢复服务的正常运行，提高了系统的可靠性和可用性。分布式架构还可以通过数据冗余、备份等手段来保障数据的安全性和完整性，进一步增强了系统的容错能力。

\paragraph{典型应用}

微服务架构是分布式架构的一种典型应用，它将系统分解为多个小型的、独立的微服务，每个微服务专注于特定的业务功能。在智慧水利平台中，可将监测接入、设备管理、告警分析、工单流转、三维场景服务等模块封装为独立微服务，实现系统的高可扩展性、灵活性和可靠性。微服务架构使得不同团队可以围绕各自服务独立开发、部署和扩展，加快平台能力迭代。

\paragraph{水利工程安全监测平台的远程监控功能}监测数据通过分布式服务处理。在该平台的远程监控场景中，分布式架构也被有效应用。监测采集服务负责从各类设备中获取数据，并将其传输给数据处理服务；数据处理服务对监测数据进行分析和清洗，判断是否发生异常情况；如果需要触发预警或执行其他操作，则通过网络调用相应服务来完成。通过这种方式，平台的各个功能模块可以分布在不同的服务器或节点上协同工作，提高系统性能和可扩展性。

表\ref{tab:arch-comparison}把前面讨论的几种结构并列比较。需要提醒的是，表中各行并不是互斥的选项：分层描述的是代码的逻辑组织方式，客户\slash 服务器描述的是交互与部署关系，分布式描述的是物理部署形态。一个系统完全可以同时具备三者，“分层的微服务分布式部署”就是常见组合，因此这张表的用法是按维度对照，而不是从中挑一个。

\begin{table}[htbp]
\centering
\caption{关键对比与选择建议}\label{tab:arch-comparison}
\begin{tabular}{p{2.8cm}p{2.8cm}p{4.2cm}p{4.2cm}}
\toprule
\textbf{架构类型} & \textbf{划分维度} & \textbf{适用场景} & \textbf{技术挑战} \\
\midrule
分层架构     & 逻辑结构风格 & 中小型系统，需清晰职责分离 & 层间通信优化 \\
客户/服务器   & 交互/部署风格 & 集中管理资源，瘦客户端需求  & 服务器性能瓶颈 \\
分布式架构   & 物理部署风格 & 高并发、大规模系统        & 服务治理、分布式事务处理 \\
\bottomrule
\end{tabular}
\end{table}

本章核心术语和角色应与全书保持同一口径，见表\ref{tab:ch03-terms}。术语表不仅用于写作，也用于接口、日志、数据库字段和测试用例命名；若业务方提出新称呼，应在需求评审中建立同义词映射。
\begin{table}[htbp]
\centering
\caption{第3章核心术语与角色口径}
\label{tab:ch03-terms}
\begin{tabular}{p{3.0cm}p{4.3cm}p{6.0cm}}
\toprule
类别 & 统一术语 & 使用边界 \\
\midrule
业务侧结论 & 预警 & 表示对风险趋势的业务提示 \\
系统侧对象 & 告警事件、告警服务 & 表示可追踪、可处置的系统对象 \\
岗位角色 & 值班员、专业分析员、审批人、运维员 & 分别承担监视确认、专业研判、方案审批、技术保障 \\
质量码 & valid、suspect、missing & 记录测值可信程度，不替代预警等级 \\
\bottomrule
\end{tabular}
\end{table}

模块边界还需要可度量的证据。扇入表示依赖某模块的调用方数量，扇出表示该模块直接依赖的模块数量；接口数量反映边界暴露面，版本联动次数反映一次变更需要同步修改多少调用方，共同变更频率则可从版本库提交记录统计。扇出过高且共同变更频繁，说明职责可能没有拆开；扇入很高的稳定模块应优先保持向后兼容。微服务划分应依次检查业务能力边界、限界上下文、变更频率差异和团队所有权：只有能独立发布、独立观测并由明确团队负责的边界，才值得成为服务；仅为追求“服务数量”而拆分会放大网络与运维成本。

在软件架构设计领域，架构风格与模式为解决常见的设计问题提供了通用的解决方案，影响着软件系统的结构、行为和质量属性。本部分讨论 MVC、微服务、事件驱动等常见的架构风格与模式。

本章的架构决策需要和后续实现章节衔接：微服务的 Spring Boot 接口与安全边界在第5章展开，前端原型与 MVC 的视图实现见第4章，三维场景订阅与空间对象服务见第6、7章，ADR 的预警等级和质量码落地见第8章。阅读这些章节时，应回到本章的职责、依赖方向和性能基线核对实现是否改变了原决策。

\subsection{MVC架构模式}

\paragraph{定义与架构组成}

MVC 是一种将软件应用程序分为三个主要部分的架构风格：模型（Model）、视图（View）和控制器（Controller）。模型代表应用程序的数据和业务逻辑，负责处理数据的存储、检索和更新；视图负责呈现数据给用户，它是用户与系统交互的界面；控制器则充当模型和视图之间的桥梁，接收用户输入，根据输入调用模型的方法进行处理，并选择合适的视图来显示结果。

以水利工程安全监测平台为例，监测数据处理与预警判定逻辑可归为模型部分，用户通过监测大屏或 Web 页面查看工程状态的界面属于视图，而处理查询条件、图层切换和告警确认等操作的部分则属于控制器。这种分离使得各部分职责明确，开发者可以专注于特定部分的开发与维护。

\paragraph{优势与代价}

MVC 的主要优势在于其高可维护性和可扩展性。由于各部分相互独立，修改其中一个部分不会对其他部分产生直接影响。例如，如果需要更新水利工程安全监测平台的监测总览界面（视图），只需修改视图部分的代码，而无需改动模型和控制器。MVC 模式还支持并行开发，不同的开发团队可以分别负责模型、视图和控制器的开发，提高开发效率。

该模式适用于需要稳定用户交互的 Web 应用；本书技术基线下可采用 Spring MVC 组织后端控制器、领域服务和视图接口。代价是控制器容易不断堆叠流程而形成“胖控制器”，模型可能退化为只保存字段的“贫血模型”，视图与模型字段的直接绑定也会造成耦合。应通过用例服务、DTO 和校验器分担职责，并用接口测试约束视图与领域模型的边界。

\paragraph{案例分析}

以一个水库安全监测平台为例，测站信息、监测记录和告警规则等属于模型部分；用户看到的监测大屏、趋势图页面和测站详情页等是视图；而处理查询条件、图层切换、告警确认等操作的部分则是控制器。通过 MVC 模式，系统的结构清晰，开发和维护变得更加高效（见图\ref{fig:mvc-architecture}）。

\begin{figure}[H]
\centering
\begin{tikzpicture}[>=Stealth, scale=0.9, thick]
  \node[draw=Primary, rectangle, minimum width=1.3cm, minimum height=0.8cm, fill=PrimaryLight, text=PrimaryDark, rounded corners=2pt, thick] at (0,1.2) (model) {\small\bfseries Model};
  \node[draw=Primary, rectangle, minimum width=1.3cm, minimum height=0.8cm, fill=PrimaryLight, text=PrimaryDark, rounded corners=2pt, thick] at (-2,0) (view) {\small\bfseries View};
  \node[draw=Primary, rectangle, minimum width=1.3cm, minimum height=0.8cm, fill=PrimaryLight, text=PrimaryDark, rounded corners=2pt, thick] at (2,0) (ctrl) {\small\bfseries Controller};
  \node[draw=Primary, circle, minimum size=0.9cm, fill=PrimaryLight, text=PrimaryDark, thick] at (-4,0) (user) {User};

  \draw[->,draw=Primary, thick,>=Stealth] (user.south) to[bend right=28] node[below,font=\scriptsize] {请求} (ctrl.south);
  \draw[->,draw=Primary, thick,>=Stealth] (ctrl) -- node[right,font=\scriptsize] {调用} (model);
  \draw[->,draw=Primary, thick,>=Stealth] (model) -- node[above,font=\scriptsize] {数据} (view);
  \draw[->,draw=Primary, thick,>=Stealth] (ctrl) -- node[left,font=\scriptsize] {选择视图} (view);
  \draw[->,draw=Primary, thick,>=Stealth] (view) -- node[below,font=\scriptsize] {页面数据} (user);
\end{tikzpicture}
\caption{MVC架构模式}
\label{fig:mvc-architecture}
\end{figure}

\subsection{微服务（Microservices）}

\paragraph{定义与架构特点}

微服务架构是一种将大型应用程序拆分为多个小型、独立服务的架构风格。每个微服务都围绕特定的业务功能构建，有自己独立的数据库、业务逻辑和接口。这些服务通过轻量级的通信机制（如 RESTful API）进行交互。

与传统的单体架构相比，微服务架构具有更高的灵活性和可扩展性。在水利工程安全监测平台中，可以将监测接入、数据治理、预警分析、用户管理等功能分别拆分成独立的微服务。这样，当需要增加新的监测点类型或调整预警算法时，只需要对相应的微服务进行调整，而不会影响系统的其他部分。

\paragraph{优势与挑战}

微服务的优势明显。一方面，每个微服务都相对较小且独立，开发者可以更轻松地理解和修改其代码，系统可维护性随之提高。另一方面，微服务架构支持独立部署和扩展，可根据业务需求对特定服务进行资源调整。例如，在汛期高峰期，可单独扩展监测接入与预警分析微服务，以应对大量实时监测数据。

然而，微服务架构也带来了一些挑战。由于服务数量众多，服务治理变得复杂，包括服务发现、负载均衡、容错处理等。分布式系统中的数据一致性问题也需要额外的技术手段来解决，例如使用分布式事务管理机制。

\paragraph{案例分析}

以省级智慧水利平台为例，监测接入、告警分析、工单管理和专题分析等功能都可以通过微服务架构实现。每个微服务独立开发、部署和扩展，能够快速响应业务需求的变化，保证系统的高可用性和高性能。微服务架构的整体结构如图\ref{fig:microservices-architecture}所示。

\begin{figure}[H]
\centering
\begin{tikzpicture}[>=Stealth, scale=0.85, thick]
  % API Gateway
  \node[draw=Primary, rectangle, minimum width=5cm, minimum height=0.6cm, fill=PrimaryLight, text=PrimaryDark, rounded corners=2pt, thick] at (0,3.5) (gw) {\small\bfseries API Gateway};

  % Services
  \node[draw=Primary, rectangle, minimum width=1.4cm, minimum height=0.5cm, fill=PrimaryLight, text=PrimaryDark, rounded corners=2pt, thick] at (-2,2) (s1) {\tiny\bfseries 接入服务};
  \node[draw=Primary, rectangle, minimum width=1.4cm, minimum height=0.5cm, fill=PrimaryLight, text=PrimaryDark, rounded corners=2pt, thick] at (0,2) (s2) {\tiny\bfseries 告警服务};
  \node[draw=Primary, rectangle, minimum width=1.4cm, minimum height=0.5cm, fill=PrimaryLight, text=PrimaryDark, rounded corners=2pt, thick] at (2,2) (s3) {\tiny\bfseries 数据服务};

  % Databases
  \node[draw=Primary, rectangle, minimum width=1.4cm, minimum height=0.5cm, fill=PrimaryLight, text=PrimaryDark, rounded corners=2pt, thick] at (-2,0.8) (d1) {\tiny\bfseries 数据库1};
  \node[draw=Primary, rectangle, minimum width=1.4cm, minimum height=0.5cm, fill=PrimaryLight, text=PrimaryDark, rounded corners=2pt, thick] at (0,0.8) (d2) {\tiny\bfseries 数据库2};
  \node[draw=Primary, rectangle, minimum width=1.4cm, minimum height=0.5cm, fill=PrimaryLight, text=PrimaryDark, rounded corners=2pt, thick] at (2,0.8) (d3) {\tiny\bfseries 数据库3};

  % Lines
  \draw[draw=Primary, thick] (gw) -- (s1);
  \draw[draw=Primary, thick] (gw) -- (s2);
  \draw[draw=Primary, thick] (gw) -- (s3);
  \draw[draw=Primary, thick] (s1) -- (d1);
  \draw[draw=Primary, thick] (s2) -- (d2);
  \draw[draw=Primary, thick] (s3) -- (d3);
\end{tikzpicture}
\caption{微服务架构示意图}
\label{fig:microservices-architecture}
\end{figure}

\subsection{事件驱动（Event-Driven）}

\paragraph{定义与工作原理}

事件驱动架构是基于事件来触发和处理操作的架构风格。系统中的模块通过发布和订阅事件进行通信，当一个模块发布特定事件时，对该事件感兴趣的其他模块会收到通知并执行相应的操作。

在水利工程安全监测平台中，当位移、渗流或库水位等监测指标超过阈值时，系统会发布异常事件。负责预警的模块订阅这类事件，当收到事件通知后，会立即触发告警动作，如生成告警工单、向值班员推送通知或联动大屏高亮显示。这种架构风格使得系统能够对外部事件做出实时响应，提高平台的灵活性和可扩展性。新的功能模块只需订阅相关事件，就可以轻松集成到系统中，而不需要对现有模块进行大规模修改。

\paragraph{优势与代价}

事件驱动架构的主要优势在于其高度的异步性和松耦合特性。异步处理使得系统能够在不阻塞其他操作的情况下处理事件，提高了系统的响应速度和吞吐量。代价是跨服务状态通常只能最终一致，事件模式演化需要兼容旧消费者，链路调试和追踪比同步调用困难，还必须设计幂等键、重试和死信处理。松耦合带来的收益只有在这些运维机制可观测、可演练时才成立。

这种架构风格适用于实时监测、告警通知和工单联动等场景。大量测点持续产生观测事件时，发布者无需等待所有下游处理完成，订阅者也可按职责独立扩展。

\paragraph{案例分析}

以大坝渗压超阈值为例，数据处理服务发布包含工程编码、测点编号、发生时间、当前值、质量码和规则版本的异常事件。告警服务生成告警并推送值班终端，工单服务创建处置任务，三维场景订阅同一事件后高亮对应坝段。各订阅者以事件标识去重；若某个订阅者暂时失败，可依据消息持久化和重试策略恢复，不影响原始数据继续入库。

\subsection{架构风格与模式的选择}

架构风格的选择直接影响部署成本、故障边界和演化方式，需要综合考虑项目规模、业务复杂度、质量目标以及团队的交付和运维能力。

架构选择应分别回答三个问题：系统分解层面采用单体、模块化单体还是微服务；每个模块内部采用 MVC 等何种组织方式；跨模块通信采用同步调用还是事件驱动。小型项目可以是“模块化单体 + MVC + 同步调用”，而汛期需要独立扩展告警分析的系统可以是“微服务 + 各服务内部 MVC + 事件驱动通知”。三项决策分别记录边界、团队能力和质量目标，不能把它们当作互斥选项。

组合方案也要显式记录代价：微服务增加网络调用、部署和分布式数据一致性成本；MVC 需要控制器、模型和视图的边界治理；事件驱动增加消息积压、重复消费和追踪成本。评审时把新增成本与 P95、可靠性、团队运维能力逐项对照，只有能提供监控、回放和故障演练证据时才接受组合。

\subsection{单体、模块化单体与微服务：用案例水库的数字做一次决策}

三种分解方式的取舍不能靠偏好，要靠本项目的数字。先把案例水库的负载算出来：28个测点按5分钟采样，稳态约每分钟5.6条观测；汛期加密到1分钟并叠加补传，峰值按每秒10条估计；并发用户是值班室2块大屏、十余个Web会话和移动巡检，合计不超过50个活跃连接；数据年增量约300万行观测，PostgreSQL单实例毫无压力。拿这组数字对照三个方案。

纯单体：一个Spring Boot进程装下全部功能。开发调试最简单，但告警评估、报表统计与数据接入共享一个进程的线程池与发布周期——报表慢查询可能拖垮告警时延，改一行展示代码也要整体重发。对本案例的负载，这些风险出现的概率低，但违反了REQ-MON-03对告警链路的独立时延承诺，属于“现在没事、出事没退路”的结构。

模块化单体：仍是一个可部署单元，但接入、质量、预警、查询按包边界严格隔离，模块间只经接口调用，事件经进程内总线传递。它保留了单体的运维简单性，又为将来拆分预留了切割线；配合消息队列做异步削峰，告警链路的时延可以单独度量与守护。按上面的数字，这是本案例的正确答案，也是配套仓库采用的结构。

微服务：接入、预警、查询各自独立部署。它换来的独立扩缩容与独立发布，在每秒10条观测、50个并发的规模下用不上；换出去的却是实打实的成本——服务间网络调用、分布式追踪、每服务一套配置与告警、跨服务数据一致性。只有当平台要接入数十座工程、团队扩到多个小组各管一域时，才值得沿模块化单体预留的边界拆分。

这个结论写成 ADR 应包含重新评估条件：接入工程数超过10座、峰值观测超过每秒200条、或出现两个以上独立交付团队时，重启本决策。决策的价值不在“选对”，而在把“什么时候要改”写清楚——第1章讲过的能力叠加逻辑在架构上同样成立。

\subsection{Web 形态的特殊考虑}

案例水库平台以浏览器为主要入口，Web 形态给架构增加三个约束，处理思路在后续章节都有落点。其一是网络密集性：页面与数据都经网络往返，静态资源应交给 CDN 或反向代理缓存并压缩传输，观测数据接口则按时间窗和聚合粒度控制单次载荷（第7章的降采样正是为此服务）。其二是并发访问：汛期大屏与移动端同时涌入时，负载均衡把请求分散到多个无状态实例，测站基础信息等热点读取放入分布式缓存，写入路径靠消息队列削峰（第5章的 Kafka 与缓存一致性详述其代价）。其三是性能优化贯穿前后端：前端做代码拆分与路由级按需加载，后端按真实查询条件设计联合索引并用执行计划验证（第5章5.7节）。这三个约束的共同点是：都不能靠单点调优解决，必须在架构分层时就预留水平扩展与缓存的位置。

Web 应用还有导航与内容组织问题：监测总览、测点详情、告警列表和处置工单构成用户的主要动线，路由结构应与业务对象一一对应（\texttt{/monitoring}、\texttt{/assets/\{id\}}、\texttt{/warnings}），使 URL 可复制、可回退、可审计；帮助文档与预案文本等内容型页面按版本管理，发布走与代码相同的评审流程，不在生产环境直接编辑。

\section{架构设计过程}

\subsection{系统环境表示}

系统环境表示是架构设计的起点：先画清系统与外部世界的边界，再谈内部分解。边界上的交互归为四类接口。用户接口面向值班员、专业分析员、审批人和运维员，形态包括监测大屏、Web 页面与移动端；硬件接口面向渗压计、GNSS 站、雨量站等设备及其遥测终端，协议与报文格式由第1章的通信规约约束；软件接口面向上级监管平台、短信网关、气象服务等外部系统，通常以 REST API 或消息队列对接；通信接口约定网络链路、带宽与安全传输要求。每类接口都要写明方向、协议、频率与失败处理，否则边界就只是一条画在图上的虚线。

约束条件与接口同样属于环境的一部分。技术约束来自选定的技术基线（本书为 Vue 3.4+/Spring Boot 3.2+/PostgreSQL 等）与既有系统的兼容要求；业务约束来自规程与权限——例如告警确认必须留痕、汛期采样频率与阈值随工况调整、只有授权角色才能修改安全参数；成本与合规约束决定了哪些“更好”的方案在本项目不可行。环境表示的交付物是一张上下文图加一张接口与约束清单，第2章的上下文图画法在此直接复用；案例水库的完整环境表示见3.6节的案例研究。

\subsection{原型定义与细化}

在软件架构设计过程中，原型定义与细化是验证架构可行性、降低项目风险的重要环节。快速原型能够让架构师和项目团队在早期阶段对架构设计进行实际检验，及时发现潜在问题并加以调整，确保最终的软件系统能够满足业务需求。

\paragraph{原型的定义与作用}

原型是软件系统的初步版本，它具有真实系统的部分功能和特性，但通常在完整性和性能上不如最终产品。对于水利工程安全监测平台而言，原型可以是一个简化版本，具备基本的监测数据接入、简单的预警功能以及基础的用户交互界面。其主要作用在于：

\paragraph{直观呈现设计理念}通过原型，项目团队能够将抽象的架构设计转化为具体可操作的实体，让利益相关者（如业务部门、开发团队成员、管理层等）直观地感受到系统的功能和使用方式。这有助于各方更好地理解架构设计，发现潜在的需求不一致或设计缺陷。例如，值班员可以通过操作原型，提前发现某些功能不符合实际使用习惯，及时提出修改意见，避免在开发后期进行大规模返工。

\paragraph{验证架构的可行性}在水利工程安全监测平台的开发中，通过构建原型，可以验证系统架构在技术上是否可行。比如，测试在模拟多个监测点同时工作的情况下，系统的数据处理和传输是否能够满足实时性要求；验证系统各模块之间的通信和协作是否顺畅，是否存在接口不兼容或交互异常的问题。如果在原型阶段发现问题，可以及时调整架构设计，降低后期开发风险和成本。

\paragraph{促进沟通与协作}原型为项目团队成员、业务部门以及其他利益相关者提供了一个共同的沟通平台。在讨论和评估原型的过程中，各方可以充分交流意见和想法，明确项目的目标和需求。开发团队可以根据业务反馈，进一步细化和优化架构设计，提高项目成功率。

\paragraph{快速原型的构建方法}

快速原型的构建需要采用一些高效的方法和工具，以快速实现原型并进行验证。

\paragraph{选择合适的技术栈}根据项目需求和团队能力，选择能够快速开发的技术框架和工具。例如，对于水利工程安全监测平台的前端原型，可以使用 HTML、CSS 和 JavaScript，结合 Vue.js 等框架快速搭建监测总览界面；后端可以选择轻量级框架快速实现数据处理和接口功能。利用模拟数据生成工具为原型提供监测数据和业务数据，以便进行功能测试。

\paragraph{采用迭代式开发}快速原型通常采用迭代式的开发方式。在每个迭代周期中，逐步增加原型的功能和完善其细节。例如，在第一个迭代中，实现基本的传感器数据显示功能；在第二个迭代中，添加预警功能的初步实现；在后续迭代中，不断优化用户界面和系统性能。通过迭代式开发，可以及时根据反馈调整原型，使其逐渐接近最终产品的要求。

\paragraph{注重核心功能的实现}在构建原型时，应重点关注系统的核心功能和业务流程。对于水利工程安全监测平台来说，监测数据接入、异常预警和用户查询等是核心功能。优先实现这些功能可以快速验证架构的可行性。对于一些辅助功能和细节，可以在后续优化阶段完善，避免在原型阶段花费过多时间和精力。

\paragraph{原型的细化与优化}

在完成初步的原型构建后，需要根据反馈对原型进行细化和优化。

\paragraph{收集反馈意见}通过组织利益相关者对原型进行试用和评估，收集各方的反馈意见。这些意见可能包括功能需求的调整、用户界面的改进建议、性能优化的要求等。例如，用户可能认为原型的操作流程不够便捷，需要简化；开发团队可能发现某些功能在技术实现上存在困难，需要调整架构设计。

在收集反馈时，应采用多样化的方法，如用户测试、问卷调查、小组讨论等，以确保获取全面且有价值的信息。对于水利工程安全监测平台，可以组织值班员和专业分析员使用原型，观察他们在监测查询、告警确认和图表分析过程中的操作行为和遇到的问题。向利益相关者发放详细问卷，涵盖功能完整性、易用性和性能表现等多个方面，并组织开发团队、产品负责人和业务代表进行讨论，共同分析问题产生的原因和可能的解决方案。

\paragraph{功能细化与扩展}根据反馈意见，对原型的功能进行细化和扩展。比如，在水利工程安全监测平台原型中，最初的预警功能可能只是简单弹窗提示，在细化阶段，可以增加告警推送功能，将预警信息及时发送到值班员的终端，并提供监测点位置、指标类型和趋势变化等信息。

进一步扩展水利工程安全监测平台的预警功能，除了向值班员终端推送预警信息外，还可以与调度平台或应急值守系统进行联动。当检测到渗流异常或位移超限时，通知值班员，也能自动将相关信息发送到调度中心，包括工程位置、异常指标、发生时间等详细信息，以便快速响应。为预警功能增加智能分类能力，根据监测数据和预设规则，区分渗流异常、变形异常、水位异常等不同类型，针对不同类型的预警提供差异化处置提示。

\paragraph{性能优化}对原型的性能进行优化，确保在实际使用场景下能够满足性能要求。例如，优化传感器数据的处理算法，提高数据处理的速度和准确性；优化系统的网络通信机制，减少数据传输的延迟。

在优化监测数据处理算法方面，可以采用更高效的数据过滤和分析算法。比如，使用机器学习算法对监测数据进行实时分析，提前识别可能出现的异常趋势，而不只是在异常发生后处理。这样可以提高系统响应速度，减少误报率。对网络通信机制进行优化，采用数据压缩和异步传输技术减少延迟；对存储层采用分布式时序存储方案，提高监测数据和视频数据的读写效率；在服务器端使用负载均衡技术应对高并发访问。通过这些性能优化措施，使水利工程安全监测平台原型在实际使用中更加稳定、高效。

\subsection{构件级设计}

构件级设计把系统架构细化为可实现、可测试的单元，主要涵盖接口、类以及构件间协作；这些设计决定职责能否独立验证、接口能否稳定演化。

\paragraph{接口设计}

接口是构件与外部交互的边界，它定义了构件提供的服务以及与其他构件进行通信的方式。良好的接口设计能够提高构件的独立性和可复用性，降低系统的耦合度。

\paragraph{接口设计原则}接口设计应遵循若干重要原则。一致性原则要求同一系统内相似功能的接口在设计上保持一致，包括接口命名规范、参数类型和顺序等。例如，在水利工程安全监测平台中，各类监测设备的数据获取接口命名统一为 readMonitoringData()，方便开发者理解和使用。最小化原则要求接口只暴露必要操作，避免功能冗余。以平台的预警接口为例，只需要暴露触发预警、确认预警等核心操作，而不需要将内部判定逻辑暴露给外部。稳定性原则要求接口一旦确定，应尽量保持稳定，避免频繁修改。

\paragraph{接口类型与实现}接口的类型多种多样，常见的有函数接口、面向对象接口以及基于消息的接口（如 RESTful API）。在水利工程安全监测平台中，若采用面向对象编程，监测设备与数据处理模块之间可以通过定义接口来规范交互方式。例如，定义一个 MonitoringDataProvider 接口，其中包含 readMonitoringData() 方法，具体的监测设备类实现该接口，这样数据处理模块就可以通过统一方法获取监测数据，而无需关心具体设备实现细节。如果涉及与外部系统交互，则可采用 RESTful API 提供获取工程状态、接收控制指令等接口服务。

\paragraph{类设计}

类是面向对象编程中的基本单元，类设计直接影响到软件的结构和功能实现。在类设计过程中，需要充分考虑类的职责、属性和方法。

\paragraph{类的职责单一性}每个类应具有明确单一的职责，这符合“单一职责原则”。例如，在水利工程安全监测平台中，AlarmManager 类只负责告警事件的管理工作，如根据预警等级生成告警事件、记录事件日志和同步事件状态，而不涉及监测数据采集或用户界面处理。这样的设计使类的功能清晰，易于理解和维护。

属性与方法的设计：类的属性用于描述类所代表的实体的特征，方法则用于实现类的行为。在设计属性时，要确保其能够准确反映类的状态，并且具有合适的访问权限。例如，Sensor类可提供 readLatestValue() 读取最新测值和 selfCheck() 执行设备自检，状态字段通过公共方法读取并校验。事件判定属于异常分析服务，传感器不直接创建预警事件，从而保持采集职责单一。

\paragraph{类的继承与多态}合理利用类的继承和多态特性可以提高代码复用性和扩展性。在水利工程安全监测平台中，可以定义一个 Sensor 基类，包含设备基本信息和数据采集的通用方法。然后，PiezometerSensor、DisplacementSensor、RainGaugeSensor 等具体设备类继承自 Sensor 基类，并根据自身特点重写或扩展相关方法。通过这种方式可以减少代码重复编写，也使平台在添加新设备类型时更容易扩展。

\paragraph{构件间的协作}

构件间的协作是实现系统整体功能的核心，它决定了各个构件如何协同工作以完成复杂的业务流程。在水利工程安全监测平台中，清晰且高效的协作机制能保障系统在监测、预警和处置场景下稳定运行。

从数据流转的角度来看，不同构件在数据的产生、传输、处理及存储环节紧密配合。传感器作为数据的源头，时刻监测着工程和环境信息，如渗压计感知坝体渗压变化、雨量计记录降雨过程、GNSS 位移监测点采集位移数据。一旦有数据产生，传感器便会将其发送给数据采集与预处理构件。该构件负责对原始数据进行初步清洗，剔除噪声和异常值，以确保数据的准确性和可用性，随后将处理后的数据传递至数据分析构件。

异常分析构件承担着重要任务，它依据预设的规则和算法，对数据进行深度挖掘和分析。例如，通过对一段时间内渗压、位移和雨量变化趋势的分析，判断工程是否存在异常风险；根据多源监测指标的组合变化，预测预警等级。若分析结果满足规则，异常分析构件向预警触发构件发布事件。预警触发构件生成告警事件，通过向值班终端推送通知等方式提示风险；事件、规则版本和处置回执同时写入日志记录构件，便于后续查询和追溯。

从系统控制层面看，值班员可通过Web端提交采集频率调整、测点停用或告警临时屏蔽请求。用户操作构件校验身份、权限和参数后，把命令交给配置服务；配置服务记录版本与有效期，再由监测接入模块向指定设备下发。告警屏蔽只影响通知策略，不删除原始事件，并须设置到期时间和复核人。设备返回执行回执后，平台更新配置状态并保留完整审计链。

当构件生产速率与消费速率不一致时，可用消息队列提供缓冲和异步通信。水利工程安全监测平台出现瞬时上传洪峰时，持久化队列允许数据处理构件按自身能力消费，避免接入线程被下游阻塞；代价是必须处理重复消息、积压、顺序和重放。服务注册与发现用于动态定位服务实例，但仍需配合超时、熔断和健康检查，不能替代接口契约与故障处理。

\section{案例研究：水利工程安全监测平台架构}

\subsection{需求分析和用例建模}

需求分析与用例建模是系统架构设计的第一步，其目标是准确理解水利业务需求，明确平台需要支撑的监测、分析、预警和处置能力。

\paragraph{需求分析}

\paragraph{功能需求}

水利工程安全监测平台的核心功能围绕“感知接入 - 数据治理 - 风险研判 - 预警处置”展开。首先，平台需要接入渗压计、测缝计、位移计、GNSS 监测站、雨量站、库水位计和视频监控等多源设备，实现对大坝变形、渗流、降雨、水位和现场工况的实时采集。其次，平台需要对多源监测数据进行统一汇聚、质量校验、时间对齐和异常剔除，形成可供分析使用的标准化数据资源。再次，平台应支持阈值预警、趋势分析、模型研判和告警联动，当监测指标持续偏离正常区间时，能够自动触发预警并推送处置建议。

平台还应具备多终端协同能力。值班员希望通过监控大屏实时掌握水库运行态势，通过 Web 端查看分测点趋势曲线和告警记录，通过移动端在巡检现场快速核对设备状态和处理结果。专业分析员需要对预警阈值、分析模型和应急流程进行研判，运维员则要完成设备注册、通信链路维护、账号权限控制以及与上级监管平台的数据交换。

\paragraph{非功能需求}

系统的可靠性和实时性是水利业务场景中的首要非功能需求。监测平台需要在汛期、强降雨、突发险情等高压工况下保持持续运行，即使出现局部断电、边缘站点离线或通信网络抖动，也要通过断点续传、缓存补报、双链路通信和主备切换等机制保证重要数据不丢失、预警链路不中断。例如，前置采集网关可在通信恢复后自动补传数据，平台中心可对关键告警采用短信、应用消息和声光告警等多通道并发推送。

平台承载工程安全数据和调度决策信息，泄露、篡改或越权操作可能直接影响研判与处置，因此需要在采集链路、平台服务、用户访问和审计追踪等层面建立防护机制。数据传输过程中应采用 TLS、VPN 或专网加密，重要操作应保留审计日志，重要告警应支持签名和确认回执，用户认证应结合统一身份认证和分级授权策略，确保监测数据、预警结论和处置指令不会被篡改或误用。

平台的易用性与可扩展性也需要一并考虑。监测界面应突出重点测点、重要告警和风险趋势，避免在紧急场景下出现信息淹没；系统架构则应支持后续接入新的监测设备、新增分析模型或扩展到河道、泵站、灌区等更多水利场景。

\paragraph{用例建模}

\paragraph{确定参与者}

本案例沿用第2章的追踪矩阵编号规则，把水库主线需求登记为 REQ-MON 系列，例如：REQ-MON-01（观测入库需带质量码与事件时间）、REQ-MON-02（黄色及以上预警须人工确认后才生成工单）、REQ-MON-03（交互查询 P95 不超过表\ref{tab:performance-baseline}基线）、REQ-MON-04（每条处置须能回查原始读数与规则版本）。后文的架构决策与第8章的阶段验收都按这些编号回指需求，读者据此可以检验“需求—架构—验收”链条是否闭合。水利工程安全监测平台的主要参与者包括值班员、专业分析员、审批人、运维员和外部协同系统。值班员负责实时查看监测总览、确认预警事件、跟踪处置状态；专业分析员负责阈值设置、分析模型维护和风险研判；审批人负责处置方案审批；运维员负责设备接入、通信维护、权限控制和运行保障；外部协同系统则包括短信网关、视频平台、气象服务和上级监管平台等，为平台提供消息通道、外部数据或监管接口。

\paragraph{识别用例}

从值班员角度来看，“实时监测总览”是最核心的用例。值班员进入平台后，需要查看水库基础信息、监测测点布局、实时告警列表和风险趋势曲线，并能快速定位异常测点。对于发生超阈值或模型判定异常的测点，系统应自动高亮显示，并提供测点最近数据、空间位置、关联视频和历史同类事件记录。

“预警处置”是另一个重要用例。当系统检测到渗压升高、位移突变、雨量骤增或水位异常波动时，应自动生成预警事件，并将事件等级、触发规则、影响测点、建议处置动作和责任人推送给值班员。值班员需要在平台中完成确认、转办、现场核查、闭环处理等动作，系统应记录每一步处理时间和责任归属。

从专业分析员和运维员角度，“设备接入配置”、“阈值与模型管理”、“用户权限管理”也是重要用例。“设备接入配置”由运维员执行，涉及新增测点、配置采集频率、绑定通信协议、维护设备台账；“阈值与模型管理”由专业分析员负责，覆盖不同工程、不同测点类型的预警规则维护；“用户权限管理”确保值班、分析、审批、运维等角色只拥有与职责匹配的操作权限。

通过对这些用例的详细分析，可以清晰展示水利工程安全监测平台的业务流程、角色分工与信息流转关系，为后续的架构分层、服务拆分和接口设计提供依据。

\subsection{构件协作与接口设计}

在水利工程安全监测平台架构中，构件协作与接口设计决定了平台是否能够稳定支撑“采、传、存、算、用”全链路业务。合理规划构件间协作关系并定义标准接口，能够确保各个模块协同工作，完成监测感知、分析预警和业务处置任务。

\paragraph{构件协作}

\paragraph{（1）传感器构件与数据处理构件协作}

监测设备构件是平台感知工程状态的“触角”，负责采集渗压、位移、变形、雨量、水位和视频巡检等数据。前置采集网关对接各类仪器后，将原始数据上传至接入服务。数据处理构件接收到监测数据后，首先进行协议解析、质量校验、异常值识别和时间对齐，再依据阈值规则、趋势模型或多指标联合分析策略进行研判。如果位移速率连续升高且同期雨量快速增加，数据处理构件应及时生成风险事件，并将结果传递给预警构件与可视化构件。

\paragraph{（2）预警触发构件与用户交互构件协作}

一旦预警构件接收到来自数据处理构件的事件指令，便立即启动预警处置流程。它需要在监控大屏和 Web 端展示告警信息，也应通过消息推送、短信或语音外呼等方式通知责任岗位。用户交互构件负责向值班员呈现预警等级、异常指标、空间位置、处置建议和历史关联事件。值班员收到通知后，可以在平台中完成确认、派单、升级上报和关闭事件等操作，必要时还可以联动视频构件查看现场画面，或调用应急预案服务生成处置清单。

\paragraph{（3）数据存储构件与其他构件协作}

数据存储构件承担着监测数据、设备台账、告警记录、巡检记录和用户操作日志的持久化任务。数据处理构件完成数据治理和风险分析后，会将结构化结果写入时序数据库、关系数据库或对象存储中，以支撑历史回溯、趋势分析和报表统计。当用户交互构件需要查询某一测点在特定时段内的变化曲线，或需要检索某次预警事件的闭环处置记录时，数据存储构件应提供稳定高效的查询能力。数据存储构件还需支持冷热分层存储、备份恢复和归档管理，满足水利行业对长期留痕和审计追踪的要求。

\paragraph{接口设计}

\paragraph{传感器接口设计}

监测设备接口是感知层与平台层交互的桥梁，必须定义统一的数据传输格式、设备标识规则和通信协议适配方式。对于现场监测设备，接口通常需要兼容 Modbus、RS485、MQTT、HTTP 或厂商私有协议，并在接入层完成协议适配与标准化封装。接口设计中应包含设备注册、测点映射、数据读取、时间同步和异常上报等能力。这样在新增测缝计、GNSS 站或视频设备时，只需扩展适配器而无需整体重构平台。

\paragraph{预警接口设计}

预警接口主要用于分析服务、消息服务和用户交互服务之间的协同。分析服务输出预警事件时，接口应定义统一的事件格式，包括工程编码、测点编号、异常指标、预警等级、触发规则、建议动作和发生时间等字段。接口既要支持同步查询，也要支持基于消息队列的异步分发，驱动消息通知、工单流转和大屏展示。预警接口还应包含状态回写机制，用于接收值班员的确认结果、处置意见和升级操作，保证事件生命周期可追踪。

\paragraph{数据存储接口设计}

数据存储接口负责协调数据服务与业务服务之间的数据交互。对于监测数据写入，接口应明确测点编码、时间戳、单位、质量标记和值来源等字段规范；对于查询场景，接口应支持按工程、测点、时间范围、数据类型和预警等级等条件进行组合检索。为了兼顾性能与可维护性，可以将时序数据查询、告警事件检索、统计报表生成和归档数据访问拆分为不同接口族，并配套分页、缓存和权限校验机制，确保多角色、多终端访问时仍能保持稳定响应。

\paragraph{架构类型对照}
表\ref{tab:arch-comparison}把分层、客户/服务器和分布式架构放在同一选择矩阵中，读者应结合并发规模、部署边界、团队能力和故障隔离要求解释每一列，而不是把三种架构误认为互斥答案。

代码清单\ref{lst:ch03-java-interface}以最小接口展示构件边界：调用方只依赖监测数据契约，设备协议和重试细节可以留在实现层。

\subsection{把设计交给评审：一次架构评审会的组织}

案例的设计成果不能以“画完了”为终点，要经过一次正式评审才能进入实现。评审会的输入是四件套：质量属性场景清单（3.1节）、架构视图（用例图、类图、部署视角）、ADR 记录、REQ-MON 追踪表。会议按第2章需求评审的角色分工进行，但检查焦点不同：业务代表核对用例覆盖与角色权限边界是否符合规程；专业分析员核对数据流上单位、质量码与模型输入的一致性；开发代表核对接口契约能否直接落到第4、5章的技术栈；运维代表核对部署视角里的监控点与回滚路径。每条意见登记为“接受修改／解释后关闭／登记为风险”三种状态之一，禁止用“会后再说”结束任何一条。设计小组在上会前应先自查五问：每条场景有度量值吗？每个构件能说出它承载哪条需求吗？每个 ADR 写了重估条件吗？部署图上标了监控点吗？风险清单是空的吗——若是，先补上再来。自查通不过就上会，浪费的是所有评审角色的时间。评审通过的标志不是全票赞成，而是：每条 REQ-MON 需求都能在架构图上指出承载它的构件，每个 ADR 都有明确的重新评估条件，风险清单上的每一项都有验证手段与责任人。这份评审记录随版本归档，是第8章阶段验收表中“架构”一行要求的证据原件。

一个失败的评审长什么样也值得知道。某届课程小组的评审记录只有一句“与会人员一致同意通过”，附一张架构图。追问之下暴露出三件事：性能场景没有度量值，“响应快”无法验证；ADR只有结论没有备选方案，无从判断是否真的比较过；风险清单为空——不是没有风险，是没人愿意在通过前夜提出来。这份“通过”的评审在第8章验收时全数返工：压测发现告警链路超时，才补做了消息队列决策，此时前端联调已经开始，返工成本是评审当时的数倍。评审的作用不是仪式，而是把返工从实现期提前到纸面期；一场没有产生任何修改项和风险项的评审，几乎可以断定没有起作用。

\section{架构评估与ATAM迷你演练}

\subsection{评估方法}

软件架构评估是确保系统架构满足业务需求、具备良好性能与可维护性的重要环节，采用科学有效的评估方法能及时发现架构设计中的潜在问题并加以优化。常见的评估方法丰富多样，各自适用于不同场景。

\paragraph{ATAM（架构权衡分析方法）}

ATAM 是一种架构评估方法，着重分析多个质量属性之间的权衡。首先把可靠性、安全性和响应时间等目标改写成带刺激、环境、响应和度量的场景；其次说明构件、连接和部署决策如何影响场景；最后识别敏感点、权衡点与风险，并为高风险假设安排验证。若设备故障场景暴露恢复时间不达标，就应修改冗余或恢复策略并重新测量，而不能只给出“提高可靠性”的定性结论。

ATAM 的效用树把质量属性按业务重要性和实现难度排序，帮助团队先验证最可能影响工程安全的场景。图\ref{fig:atam-utility-tree}以“预警处置”作为根节点，向下展开可靠性、性能和安全性，再落到可测的叶节点；叶节点的优先级应在需求评审中确认，不能由架构师单方面决定。
\begin{figure}[H]
\centering
\begin{tikzpicture}[node distance=8mm and 13mm, every node/.style={font=\small}, box/.style={draw=Primary,rounded corners,fill=PrimaryLight,align=center,minimum width=2.4cm,minimum height=.55cm}, leaf/.style={draw=TextGray,fill=white,align=center,minimum width=3.0cm,minimum height=.48cm}]
\node[box] (root) {预警处置质量};
\node[box,below left=of root] (rel) {可靠性};
\node[box,below=of root] (perf) {性能效率};
\node[box,below right=of root] (sec) {安全性};
\node[leaf,below=of rel] (rleaf) {事件不丢失\\恢复后可重试};
\node[leaf,below=of perf] (pleaf) {告警生成 $P95\leq 5$ s};
\node[leaf,below=of sec] (sleaf) {确认与处置\\全量审计};
\draw[-{Stealth}] (root)--(rel); \draw[-{Stealth}] (root)--(perf); \draw[-{Stealth}] (root)--(sec);
\draw[-{Stealth}] (rel)--(rleaf); \draw[-{Stealth}] (perf)--(pleaf); \draw[-{Stealth}] (sec)--(sleaf);
\end{tikzpicture}
\caption{ATAM 预警处置效用树}
\label{fig:atam-utility-tree}
\end{figure}

\paragraph{ATAM迷你演练}设定与具体方案无关的场景“汛期接入流量达到平时10倍时，P95告警生成时延不超过5秒；任一处理节点在处理中崩溃时，已被系统确认接收的告警事件不丢失，恢复后仍能完成推送”。候选方案A为分析服务同步调用通知服务，方案B为分析服务写入持久化消息队列、通知服务独立消费。评估结果见表\ref{tab:atam-mini}。

\begin{table}[htbp]
\centering
\caption{告警链路ATAM迷你演练}\label{tab:atam-mini}
\begin{tabular}{p{2.2cm}p{4.7cm}p{5.0cm}}
\toprule
评估项 & 方案A：同步调用 & 方案B：持久化消息队列 \\
\midrule
性能 & 空载链路短；通知服务变慢会占用分析线程 & 增加入队开销；可通过分区和消费者扩容稳定吞吐 \\
可靠性 & 下游故障可能使告警生成失败 & 确认写入后可重试，但需处理重复消费 \\
复杂度 & 部署与追踪较简单 & 增加消息代理、幂等键、积压监控和回放机制 \\
决策 & 不满足单节点故障场景 & 采用方案B；以事件号作幂等键，并设积压告警阈值 \\
\bottomrule
\end{tabular}
\end{table}

这个演练把“性能更好”转化为可测场景，并明确可靠性收益以运维复杂度为代价。ATAM 交付物还应包含方法文档、效用树、风险点、非风险点、敏感点和权衡点：例如，重复消费与积压是风险点，消息分区数和消费者并发数是敏感点，已经通过故障注入和10倍流量压测验证的事件持久化可列为非风险点。评审结论应把这些判断、验证证据和后续任务一并归档。

\paragraph{SAAM（软件架构分析方法）}

SAAM 主要关注架构对系统需求变更的适应性。第一步同样是收集系统需求，明确系统的功能与非功能需求。在水利工程安全监测平台中，功能需求可能包括监测接入、预警触发和趋势分析等，非功能需求涵盖性能、可维护性等。接着，构建场景库，根据需求分析结果，创建一系列描述系统可能面临的使用场景和变更场景。比如，场景可以是增加新的监测点类型、调整预警阈值等。通过 SAAM 评估，能够提前发现架构在应对需求变更时存在的薄弱环节，增强架构的灵活性与可扩展性。

\paragraph{性能建模与仿真}

该方法通过建立数学模型来模拟系统在不同负载条件下的性能表现。对于水利工程安全监测平台，首先要确定性能指标，如响应时间、吞吐量、资源利用率等。然后，依据系统架构和预期工作负载构建性能模型。例如，利用排队论模型描述监测数据采集、传输和处理过程，考虑数据到达时间间隔、处理时间以及可能出现的排队等待情况。若仿真结果显示高负载情况下系统响应时间过长，可进一步分析瓶颈所在，并针对性提出优化措施。

\subsection{复杂度与风险分析}

在软件架构评估与验证过程中，复杂度与风险分析是重要的环节，它能帮助项目团队提前洞察潜在问题，确保架构的稳健性与项目的顺利推进。

\paragraph{软件架构复杂度分析}

\paragraph{结构复杂度}

软件架构的结构复杂度主要体现在系统的模块数量、层次结构以及模块间的依赖关系。以水利工程安全监测平台为例，如果架构包含众多相互关联的模块，如监测接入模块、预警处理模块、用户界面模块、数据存储模块等，且这些模块之间存在复杂的调用和数据传递关系，结构复杂度会增加。可以通过计算模块耦合度和内聚度来量化结构复杂度。

\paragraph{算法复杂度}

算法复杂度关乎系统处理任务时所采用算法的效率。在水利工程安全监测平台的异常识别功能中，若采用复杂的机器学习算法来分析监测数据，虽然算法可能具有较高准确性，但计算资源消耗也会增大。当大量监测数据同时涌入时，可能导致系统响应延迟。开发团队需要在算法准确性和复杂度之间找到平衡。

\paragraph{数据复杂度}

数据复杂度涉及系统中数据的规模、类型、格式以及数据之间的关系。水利工程安全监测平台可能需要处理多种类型的数据，包括监测设备采集的实时数据、工程台账、视频信息和历史预警记录等。28个测点（参数见8.1节）按5分钟一次采样，一年就是约300万条读数；测点扩到百级、采样加密到1分钟，量级还会再涨一到两个数量级，分区策略、归档周期和查询预算都必须提前设计。

\paragraph{风险分析}

\paragraph{技术风险}

技术风险通常源于采用了不成熟或不稳定的技术。在水利工程安全监测平台开发中，如果选用了尚不成熟的设备通信协议，可能会在实际应用中出现兼容性问题，导致监测设备与系统其他部分无法正常通信。应对技术风险，项目团队需要在技术选型阶段进行充分调研和验证。

\paragraph{业务风险}

业务风险与项目的业务目标和需求变化紧密相关。若水利工程安全监测平台的业务需求发生重大变化，例如新增更高等级的监测与预警要求，而现有架构难以快速支持这些新功能，就会导致项目面临业务风险。为应对业务风险，项目团队需要与业务方保持密切沟通，并在架构设计中预留一定扩展性和灵活性。

\paragraph{项目管理风险}

项目管理风险主要出在进度安排、跨专业协作和人力投入这三处。在水利工程安全监测平台项目中，如果项目进度安排不合理，导致重要任务之间的时间过于紧密，一旦某个任务延误，就可能影响整体交付时间。通过制定合理项目计划、建立有效沟通机制以及优化资源分配方案，可以降低项目管理风险。

前述评估方法在工业界的完整形态是 ATAM（Architecture Tradeoff Analysis Method）：由独立评估组主持、历时数日、产出敏感点、权衡点与风险清单的正式过程。课堂没有数日，但方法的骨架——用场景驱动走查、区分三类发现、只交清单不判优劣——完全可以在90分钟内走完，而且这恰恰是学生最缺的训练：他们见过很多架构图，几乎没见过架构被系统性质询的过程。

\subsection{课堂演练：一次90分钟的迷你ATAM}

把评估方法落成一次可以在课堂完成的演练。对象就用3.6节的案例水库架构，参加者按第2章的评审角色分工：一名学生任主持人，其余分别扮演值班员代表、专业分析员、审批人、开发与运维。演练分五步，时间按90分钟课堂切分。

第一步（10分钟）陈述架构：设计小组用一张分层图和一张部署图讲清结构，只许讲“是什么”，不许辩护“为什么好”。第二步（15分钟）确认驱动因素：从 REQ-MON 系列需求中挑出本次评估的质量属性场景，至少覆盖性能（REQ-MON-03）、可审计（REQ-MON-04）和一条可靠性场景；每条场景按“刺激—环境—响应—度量”四要素写在黑板上。第三步（30分钟）沿效用树逐场景走查：对照图\ref{fig:atam-utility-tree}，把每条场景在架构图上走一遍数据路径，记录三类发现——敏感点（一个决策影响一个属性，如消息队列容量影响告警时延）、权衡点（一个决策同时拉扯两个属性，如同步审计写入拉高可靠性却拉低吞吐）、风险（尚无证据支撑的假设，如“Kafka 不会成为瓶颈”）。第四步（20分钟）质询与记录：非设计角色每人至少提出一个失败场景（断网、汛期峰值、审批人不在岗），设计小组现场指出架构的应对路径或承认缺口；主持人把缺口登记为待验证项，标明验证手段（压测、演练或原型）与责任人。第五步（15分钟）输出结论：会议不判架构“好坏”，只交付三张清单——敏感点清单、权衡点清单、风险与待验证清单，并把它们挂到对应的 ADR 与需求编号上。

给第三步一个走查示范，学生照此粒度记录。取性能场景（REQ-MON-03）：刺激从“网关每秒推送500条观测”开始，在架构图上指第一站——接入服务。问：单实例还是多实例？答：模块化单体双实例，负载均衡分发。问：500条/秒经过质量检查的耗时？答：范围与时序校验为纯内存判断，压测单实例可达每秒2000条，不是瓶颈。下一站消息队列。问：告警评估消费滞后怎么办？答：分区数为4，评估并行度可扩——记为敏感点（分区数影响REQ-MON-03）。下一站数据库。问：批量入库与告警查询争抢连接池吗？答：写入走独立连接池，但审计写入与业务写入同库同池——记为权衡点（审计完整性对时延）。最后一站推送。问：P95从哪测到哪？答：从网关收包时间戳到值班端收到推送，埋点尚未实现——记为风险（未验证：端到端埋点方案，验证手段为联调日志比对，责任人第2组）。一条场景走完约8分钟，产出2个敏感点、1个权衡点、1个风险，全部落在清单上；没有一句“我觉得性能没问题”。

这个演练刻意保留了正式 ATAM 的骨架而砍掉了它的会议成本。学生最常见的两个失误：一是把第三步走成设计辩护会，主持人要不断把讨论拉回“这条场景在这张图上怎么走”；二是把风险写成“可能有性能问题”这类不可验证的表述，正确的写法是“未验证：单实例 Kafka 在每秒500条观测下的端到端时延，验证手段为压测脚本，责任人某某”。演练产物直接作为课程项目的架构评审证据，也是第8章阶段验收表中“架构”一行的输入。

一次真实课堂演练的产物长这样，供对照检验自己的记录粒度。敏感点清单（节选）：“消息队列分区数决定告警并行度——影响REQ-MON-03，证据为压测报告P95曲线”；“审计表与业务表同库——影响REQ-MON-04的留痕完整性，也影响备份策略”。权衡点清单（节选）：“观测入库走同步校验：提高数据质量（REQ-MON-01），但峰值时延上升约四成，缓解手段是把范围校验前移到网关”；“告警推送即时化：缩短响应时间，但重复推送风险上升，需幂等键抑制”。风险与待验证清单（节选）：“未验证：TimescaleDB连续聚合在补传乱序数据下的正确性——验证手段为注入乱序样本比对聚合值，责任人第3组，期限第12教学周”；“假设：值班员确认动作在30秒内完成——若不成立，超时升级路径的时间参数需重定，验证手段为值班室观察记录”。三张清单合计通常不超过一页，超过一页说明场景选得太散，应回到第二步重新聚焦。

\section*{小结}

架构不是一张结构图，而是一组必须写下理由的取舍。本章反复使用的做法，是把“选什么架构”拆成三个互不相同的问题：系统整体分解为单体、模块化单体还是微服务；单个模块内部采用什么方式组织（MVC 是其中一种）；跨模块通信走同步调用还是事件驱动。三个问题各有各的答案，“模块化单体 + MVC + 同步调用”和“微服务 + 服务内 MVC + 事件驱动”都是成立的组合。关键对比表要按维度对照着读，而不是从中挑一行——分层讲的是代码的逻辑组织，客户\slash 服务器讲的是交互与部署关系，分布式讲的是物理部署形态，一个系统完全可以同时具备这三样。

每一种组合都带着代价，代价必须显式写出来。微服务增加网络调用、部署和分布式数据一致性成本；事件驱动增加消息积压、重复消费和链路追踪成本。架构决策记录（ADR）就是固定这份代价的地方，它记下编号、上下文、决策、备选方案、后果、关联需求和责任人。案例水库告警链路的 ADR-003 之所以选择持久化消息队列而不是同步调用或内存队列，理由是汛期入站峰值升高且通知服务可能短时不可用；代价则是多出消息代理、幂等键、积压监控和回放运维。日后若替换实现，应新建记录并把旧记录标为“已取代”，而不是回头改写旧文档。

模块化与关注点分离、抽象与信息隐藏、功能独立性这些原则，落到案例水库的监测接入模块上都是可以当场检验的：换一种遥测终端要改动几个文件，能不能只替换适配层而不动风险评分逻辑。架构评估方法和复杂度分析给出的正是这类问题的量化答案——架构好不好，最终由这些答案说明，而不是由图画得整齐与否说明。

本轮修订把三件“怎么做”的事补进了本章：质量属性先写成带度量的场景（表\ref{tab:ch03-quality-scenarios}），单体、模块化单体与微服务用案例水库的真实负载数字做决策而不是凭偏好站队，架构评审与迷你ATAM给出了可以在一次课堂里走完的流程与产物样例（敏感点、权衡点、风险三张清单）。这三件事串起来就是本章的方法主线：场景定义“要什么”，决策记录“为什么这样选”，评审与演练检验“选得对不对”——三者共用 REQ-MON 编号，第8章的阶段验收将沿同一编号回查到今天写下的每一条场景。学生完成课程项目时，交付物不是一张架构图，而是场景表、ADR、评审记录三件套；图只是它们的插图。
\section*{章末交付物}

完成本章学习后，建议提交一份《系统架构设计说明》，用于承接第2章的需求成果。该说明至少应包含：

\begin{itemize}
\item 系统总体架构图和主要部署方式；
\item 核心模块划分，如监测接入、告警分析、数据管理、巡检工单和可视化展示；
\item 关键接口与数据交换关系；
\item 架构选型说明，包括为何采用分层、模块化单体或微服务等方案。
\end{itemize}

\section*{思考题与练习题}

\begin{enumerate}
\item 分析模块化设计原则与系统复杂度的关系，并结合依赖方向、接口数量和共同变更频率识别高耦合模块。
\item 比较MVC架构、微服务架构和事件驱动架构各自的特点，指出它们在智慧水利平台中的适用情景及优缺点。
\item 论述架构设计中的抽象与信息隐藏原则对系统可维护性和可扩展性的重要性。
\item 解释为什么需要进行架构评估与验证，列举正文介绍的评估方法，并说明如何按风险选择合适方法。
\item 对于一个具有大规模数据处理需求的水利平台，分析应该选择怎样的架构类型，并从性能、可维护性等多个角度说明理由。
\item 说明分层架构、客户/服务器架构和分布式架构分别属于哪一划分维度，并举例说明三者如何在同一水利系统中共存。
\item 论述在架构设计过程中如何平衡功能需求、非功能需求、技术成本与风险之间的关系，说明各类风险的识别与控制方法。
\item 分析Web应用架构在处理并发访问和性能优化时的特殊考虑因素，说明如何进行负载均衡和缓存优化。
\item 设计一个中等规模水利管理平台的系统架构，明确各个层级的职责、核心模块划分和模块间的接口。
\item 针对一个实际的智慧水利系统（如防洪决策支持系统），绘制用例图、类图和时序图等UML设计文档。
\item 以水利工程安全监测平台为基础，分析其架构设计中的主要决策点、权衡因素和可改进之处。
\item 设计一个水库监测系统的微服务架构，明确服务划分原则、服务间通信机制和数据管理策略。
\end{enumerate}
